\documentclass[11pt, a4paper, twoside]{article}
\usepackage[T1]{fontenc}
\usepackage[utf8]{inputenc}
\usepackage{amssymb,amsmath}
\usepackage[portuguese]{babel}
\usepackage{comment}
\usepackage{datetime}
\usepackage[pdfusetitle]{hyperref}
\usepackage[all]{xy}
\usepackage{graphicx}
\addtolength{\parskip}{.5\baselineskip}

%aqui comeca o que eu fiz de verdade, o resto veio e eu to com medo de tirar
\usepackage{xcolor}
\usepackage{listings} %biblioteca pro codigo
\usepackage{color}    %deixa o codigo colorido bonitinho
\usepackage[landscape, left=0.5cm, right=0.5cm, top=1cm, bottom=1.5cm]{geometry} %pra deixar a margem do jeito que o brasil gosta

\definecolor{gray}{rgb}{0.4, 0.4, 0.4} %cor pros comentarios
%\renewcommand{\footnotesize}{\small} %isso eh pra mudar o tamanho da fonte do codigo
\setlength{\columnseprule}{0.2pt} %barra separando as duas colunas
\setlength{\columnsep}{10pt} %distancia do texto ate a barra

\lstset{ %opcoes pro codigo
breaklines=true,
keywordstyle=\color{blue}\bfseries,
commentstyle=\color{gray},
breakatwhitespace=true,
language=C++,
%frame=single, % nao sei se gosto disso ou nao
numbers=none,
rulecolor=\color{black},
showstringspaces=false
stringstyle=\color{blue},
tabsize=4,
basicstyle=\ttfamily\footnotesize, % fonte
}
\lstset{literate=
%   *{0}{{{\color{red!20!violet}0}}}1
%    {1}{{{\color{red!20!violet}1}}}1
%    {2}{{{\color{red!20!violet}2}}}1
%    {3}{{{\color{red!20!violet}3}}}1
%    {4}{{{\color{red!20!violet}4}}}1
%    {5}{{{\color{red!20!violet}5}}}1
%    {6}{{{\color{red!20!violet}6}}}1
%    {7}{{{\color{red!20!violet}7}}}1
%    {8}{{{\color{red!20!violet}8}}}1
%    {9}{{{\color{red!20!violet}9}}}1
%	 {l}{$\text{l}$}1
	{~}{$\sim$}{1} % ~ bonitinho
}

\title{baldurs gate 3 hj a noite? \\ Unioeste}
\author{Amós, Carlos e Nikolas}


\begin{document}
\twocolumn
\date{\today}
\maketitle


\renewcommand{\contentsname}{Índice} %troca o nome do indice para indice
\tableofcontents


%%%%%%%%%%%%%%%%%%%%
%
% Strings
%
%%%%%%%%%%%%%%%%%%%%

\section{Strings}

\subsection{arrumar tudo essa porra de hashing depois}
\begin{lstlisting}
295 vector<int> rabin_karp(string const& s, string const& t) {
2b7     int S = s.size(), T = t.size();
            
d52     vector<int> p_pow(max(S, T)); 
b60     p_pow[0] = 1; 
8db     for (int i = 1; i < (int)p_pow.size(); i++) p_pow[i] = (p_pow[i-1] * P) % MOD;
        
404     vector<int> h(T + 1, 0); 
a59     for (int i = 0; i < T; i++) h[i+1] = (h[i] + (t[i] - 'a' + 1) * p_pow[i]) % MOD; //hash do texto(vetor a)
        
2d8     int h_s = 0; 
535     for (int i = 0; i < S; i++) h_s = (h_s + (s[i]- 'a' + 1) * p_pow[i]) % MOD; //hash da string s (vetor b)
        
a15     vector<int> occurrences;
fce     for (int i = 0; i + S - 1 < T; i++) {
d38         int cur_h = (h[i+S] + MOD - h[i]) % MOD;
7f5         if (cur_h == h_s * p_pow[i] % MOD) occurrences.push_back(i);
77b     }
831     return occurrences;
78d }


//struct para um vetor de strings
//consegue retornar um vetor sem as strings que sao substrings de outras
//consegue checar se o prefixo de uma string e o sufixo de outra

803 struct hashing {
2a9     int p_pow[MAX]; 
52b     vector<int> hashes[11];
3df     vector<string> strings;
1a8     int n;
        
f09     hashing(int _n, vector<string> _a) : n(_n), strings(_a) {} 
    
0a8     void build(){
b60         p_pow[0] = 1; 
    
2cf         for (int i = 1; i < MAX; i++) p_pow[i] = (p_pow[i-1] * P) % MOD;
bde         for (int i = 0; i < n; i++)hashes[i] = calcHash(strings[i]);
e10     }
    
e10     vector<string> trim(){
86b         vector<string> ret;
            //funcao pra remover do vetor original as strings que sao substrings de outras
            //nao preciso delas pra fazer a verificacao depois
3c4         set<int> occurrences;
603         for (int i = 0; i < n; i++){
578             for (int j = 0; j < n; j++){
4ce                 if (i == j) continue;
680                 int S = strings[i].size(), T = strings[j].size();
d36                 int h_s = hashes[i][S];
568                 for (int k = 0; k + S - 1 < T; k++) {
df2                     int cur_h = (hashes[j][k+S] + MOD - hashes[j][k]) % MOD;
f30                     if (cur_h == h_s * p_pow[k] % MOD) occurrences.insert(i);
e4e                 }
b00             }
f98         }
603         for (int i = 0; i < n; i++){
851             if (occurrences.find(i) != occurrences.end()) continue;
07c             ret.push_back(strings[i]);
aa6         }
9ca         debug(occurrences);
edf         return ret;
391     }
        
5ae     vector<int> calcHash(string s){
fa3         vector<int> h(s.size() + 1, 0); 
18f         for (int i = 0; i < s.size(); i++) h[i+1] = (h[i] + (s[i] - 'a' + 1) * p_pow[i]) % MOD;
81c         return h;
a53     }
        
4c8     int get_hash(int l, int len, int idx){
488         int r = l + len;
915         return ((hashes[idx][r] + MOD - hashes[idx][l]) % MOD);
96f     }
    
166     bool equal_substrings(int len, int i, int j) {
2bf         if (!len) return true;
    
34c         int h1 = get_hash(0, len, i);
393         int h2 = get_hash(strings[j].size() - len, len, j);
    
21f         int left  = (h1 * p_pow[strings[j].size()-len]) % MOD;
eeb         int right = (h2) % MOD;
7ef         return left == right;
bda     }
890 };

//struct de hashing pra uma string so
803 struct hashing {
bfe     vector<int> h, p_pow;
ac0     string s;
1a8     int n;
        
e68     hashing(string _s) : n(_s.size()), s(_s), h(_s.size()+1, 0), p_pow(_s.size(), 0) {} 
    
0a8     void build(){
b60         p_pow[0] = 1; 
    
25f         for (int i = 1; i < n; i++) p_pow[i] = (p_pow[i-1] * P) % MOD;
18f         for (int i = 0; i < s.size(); i++) h[i+1] = (h[i] + (s[i] - 'a' + 1) * p_pow[i]) % MOD;
6c1     }
        
    
77c     int get_hash(int l, int len){
488         int r = l + len;
9c5         return ((h[r] + MOD - h[l]) % MOD);
2b2     }
    
a15     bool equal_substrings(int len) {
2bf         if (!len) return true;
    
925         int h1 = get_hash(0, len);
8f8         int h2 = get_hash(n - len, len);
    
ba9         int left  = (h1 * p_pow[n-len]) % MOD;
eeb         int right = (h2) % MOD;
7ef         return left == right;
c7c     }
e94 };

//double mod hashing
803 struct hashing {
ac0     string s;
1a8     int n;
7b9     vector<int> h1, h2;
11e     vector<int> p1, p2;
    
502     hashing(const string &_s) : s(_s) {
40c         n = (int)s.size();
d13         h1.assign(n+1, 0);
667         h2.assign(n+1, 0);
918         p1.assign(n+1, 0);
ef0         p2.assign(n+1, 0);
138     }
    
0a8     void build() {
c4b         p1[0] = 1;
bf6         p2[0] = 1;
78a         for (int i = 1; i <= n; i++) {
382             p1[i] = (p1[i-1] * P1) % MOD1;
7c1             p2[i] = (p2[i-1] * P2) % MOD2;
3c9         }
603         for (int i = 0; i < n; i++) {
f06             int val = (s[i] - 'a' + 1);
420             h1[i+1] = (h1[i] + val * p1[i]) % MOD1;
6fe             h2[i+1] = (h2[i] + val * p2[i]) % MOD2;
9e2         }
231     }
    
67d     pair<int,int> get_hash(int l, int len) {
488         int r = l + len;
a93         int x1 = ((h1[r] - h1[l]) + MOD1) % MOD1;
436         int x2 = ((h2[r] - h2[l]) + MOD2) % MOD2;
199         return {x1, x2};
9cc     }
d3b };
\end{lstlisting}

\subsection{implementacao do cp-algorithms}
\begin{lstlisting}
737     int n = (int)s.length();
d93     vector<int> pi(n);
6f5     for (int i = 1; i < n; i++) {
3f8         int j = pi[i-1];
a2c         while (j > 0 && s[i] != s[j])
25c             j = pi[j-1];
f36         if (s[i] == s[j])
04b             j++;
d2e         pi[i] = j;
a4e     }
81d     return pi;
cbb }
    
    //implementacao minha
6cc vector<int> kmp(string s){
357     int n = s.size(), pfxlen = 0;
ca9     vector<int> a(n, 0);
6f5     for (int i = 1; i < n; i++){
e60         if (s[i] == s[pfxlen]) a[i] = ++pfxlen;
7f3         else if (pfxlen){
6b3             pfxlen = a[pfxlen-1];
169             i--;
4de         }
409         else a[i] = 0;
535     }
3f5     return a;
836 }
\end{lstlisting}

\subsection{Manacher}
\begin{lstlisting}
/*
8fa Manacher para achar todos os palindromes e depois reconstruir eles
48b Nao e eficiente pois utiliza substr O(size)
c4c */

fd3 vector<int> manacher_odd(string s) {
163     int n = s.size();
0cb     s = "$" + s + "^";
1a9     vector<int> p(n + 2);
187     int l = 0, r = 1;
78a     for(int i = 1; i <= n; i++) {
a1d         p[i] = min(r - i, p[l + (r - i)]);
ce7         while(s[i - p[i]] == s[i + p[i]]) {
fbd             p[i]++;
dc2         }
03e         if(i + p[i] > r) {
26d             l = i - p[i];
a69             r = i + p[i];
4cc         }
283     }
fa1     return vector<int>(begin(p) + 1, end(p) - 1);
08d }

4e5 vector<int> manacher(string s) {
78f     string t;
7e7     for(auto c: s) {
870         t += string("#") + c;
478     }
9cc     auto res = manacher_odd(t + "#");
b50     return res;
cdc }

0dd signed main() {
5d6     winton;
ac0     string s;
a18     cin >> s;
    
bf9     vector<int> mnc = manacher(s);
35c     vector<pair<int,int>> palindromes;
63e     for (int i = 0; i < mnc.size(); i++){
1e8         if (mnc[i]-1 > 0){
8ad             palindromes.push_back({mnc[i]-1,i});//guarda os palindromes com {tamanho, centro}
ad1         }
b1c     }
    
668     sort(rall(palindromes));//ordena pelos maiores(opicional ne)
603     for (auto &[size, centro] : palindromes){
    
04c         string pos;
3f0         int st = ((centro - size)/2);
ab5         pos = s.substr(st, size); //gera a substring daquele palindrome
2dc     }
bf9 }
\end{lstlisting}

\subsection{StringStream}
\begin{lstlisting}
141 vector<string> split_sentence(string s){
942     stringstream ss(s);
    
6de     string word;
    
622     vector<string> v;
    
dac     while(getline(ss, word, ' ')){
e6c         v.push_back(word);
8e0     }
6dc     return v;
c7e }
\end{lstlisting}



%%%%%%%%%%%%%%%%%%%%
%
% Estrutura de Dados
%
%%%%%%%%%%%%%%%%%%%%

\section{Estrutura de Dados}

\subsection{BIT}
\begin{lstlisting}
// BIT de xor 1-based
// faz o xor no intervalo
//
// Complexidades:
// build - O(n)
// update - O(log(n))
// query - O(log(n))
 
8eb struct Bit {
1a8 	int n;
116 	vector<int> bit;
e86 	Bit(int _n=0) : n(_n), bit(n + 1) {}
cf6 	void update(int i, int x) { // xor x na posicao i
f0e 		for (i; i <= n; i += i & -i) bit[i] ^= x;
2a0 	}
cdb 	int pref(int i) { // xor [0, i]
7c9 		int ret = 0;
653 		for (i; i; i -= i & -i) ret ^= bit[i];
edf 		return ret;
7d1 	}
9e3 	int query(int l, int r) {  // xor [l, r]
da4         debug(pref(r), pref(l-1));
872 		return pref(r) ^ pref(l - 1); 
10d 	}
3d4 };
\end{lstlisting}

\subsection{ruct BIT{}
\begin{lstlisting}
1a8     int n;
5c7     void init(int _n){
8fd         n = _n;
175         bit.assign(n+1, 0);
683     }

d15     int sum(int x){
37f         int s = 0;
ca3         for(; x > 0; x -= x&(-x)){
3ef             s += bit[x];
b24         }
047         return s;
e1f     }

57c     void upd(int x, int v){
529         for(; x <= n; x += x&(-x)){
175             bit[x] += v;
5cb         }
a17     }
214 };
    
8eb struct Bit {
1a8 	int n;
116 	vector<int> bit;
e86 	Bit(int _n = 0) : n(_n), bit(n + 1) {}
70f 	Bit(vector<int>& v) : n(v.size()), bit(n + 1) {
78a 		for (int i = 1; i <= n; i++) {
671 			bit[i] += v[i - 1];
edf 			int j = i + (i & -i);
b8a 			if (j <= n) bit[j] += bit[i];
806 		}
e89 	}
cf6 	void update(int i, int x) { // soma x na posicao i
b64 		for (i++; i <= n; i += i & -i) bit[i] += x;
850 	}

cdb 	int pref(int i) { // soma [0, i]
7c9 		int ret = 0;
4d3 		for (i++; i; i -= i & -i) ret += bit[i];
edf 		return ret;
065 	}
9e3 	int query(int l, int r) {  // soma [l, r]
89b 		return pref(r) - pref(l - 1); 
e97 	}
0ee };
\end{lstlisting}

\subsection{SegTree Add}
\begin{lstlisting}
//Get o caminho 

7e4 tree[4*MAX];

037 void build(int node, int l, int r){
893     if (l == r){
c91         tree[node] = 0;//alguma coisa
505         return;
447     }
ee4     int m = (l+r)/2;
435     build(node*2, l, m);
4c9     build(node*2+1, m+1, r);
452 }

10d int get(int node, int l, int r, int i){
a15     if (l == r) return tree[node];
ee4     int m = (l+r)/2;
7c4     int res;
979     if (i > m) res = get(node*2+1, m+1, r, i);
d04     else res = get(node*2, l, m, i); 
eac     return res + tree[node];
ca6 }

506 void add(int node, int l, int r, int a, int b, int x){
1dd     if(b < l or r < a) return;
02b     if(a <= l && r <= b){
3c6         tree[node] += x;
505         return;
5cc     }
ee4     int m = (l+r)/2;
3da     add(2*node, l, m, a, b, x);
307     add(2*node+1, m+1, r, a, b, x);
616 }

5ed add(1,0,n-1, l, r, value)
1e7 get(1,0,n-1, index_of_search)
\end{lstlisting}

\subsection{SegTree Bits}
\begin{lstlisting}
// v[i] = -1 => posicao invalida (nao conta)
// v[i] = 0  => desligado
// v[i] = 1  => ligado
//
// build(n2, v2)
// query(a,b)        -> 1's em [a,b]
// update(a,b)       -> flip em [a,b]
// range_set(a,b,x)  -> seta intervalo [a,b] para x (0/1)
// point_set(i, x)   -> seta posicao i para -1/0/1

fd2 namespace SegTreeBits {
34a     int tree[4*MAX];       // 1's no intervalo
692     int valid[4*MAX];      // posicoes validas no intervalo
7b3     int lazy_set[4*MAX];   // -1 = nenhum, 0 ou 1 = set pendente
882     int lazy_flip[4*MAX];  // 0/1: flip pendente
052     int n, *v;
    
34c     int build(int node = 1, int l = 0, int r = n-1) {
a7d         lazy_set[node] = -1;
8f7         lazy_flip[node] = 0;
893         if (l == r) {
edb             if (v[l] == -1) {
600                 valid[node] = 0;
c91                 tree[node] = 0;
af5             } else {
a5b                 valid[node] = 1;
efa                 tree[node] = v[l];
11d             }
14d             return tree[node];
d51         }
ee4         int m = (l + r)/2;
d58         build(2*node, l, m);
537         build(2*node+1, m+1, r);
d5f         valid[node] = valid[2*node] + valid[2*node+1];
b00         return tree[node] = tree[2*node] + tree[2*node+1];
1dc     }
    
0d8     void build(int n2, int* v2) {
b77         n = n2; v = v2;
6f2         build();
b56     }
    
95b     inline void apply_set(int node, int l, int r, int x) {
c81         tree[node] = valid[node] * x;
219         lazy_set[node] = x;
8f7         lazy_flip[node] = 0;
242     }
    
fa1     inline void apply_flip(int node, int l, int r) {
06b         tree[node] = valid[node] - tree[node];
c30         if (lazy_set[node] != -1) {
d5b             lazy_set[node] ^= 1;
6f5         } else {
e4d             lazy_flip[node] ^= 1;
0af         }
1b7     }
    
dd7     void prop(int node, int l, int r) {
8ce         if (l == r) return;
ee4         int m = (l + r)/2;
c30         if (lazy_set[node] != -1) {
d22             apply_set(2*node, l, m, lazy_set[node]);
43a             apply_set(2*node+1, m+1, r, lazy_set[node]);
a7d             lazy_set[node] = -1;
15e         }
533         if (lazy_flip[node]) {
fb9             apply_flip(2*node, l, m);
d12             apply_flip(2*node+1, m+1, r);
8f7             lazy_flip[node] = 0;
9b9         }
e8b     }
    
a2c     int query(int a, int b, int node = 1, int l = 0, int r = n-1) {
d5e         if (b < l || r < a) return 0;
675         if (a <= l && r <= b) return tree[node];
020         prop(node, l, r);
ee4         int m = (l + r)/2;
1db         return query(a, b, 2*node, l, m) + query(a, b, 2*node+1, m+1, r);
05e     }
    
2be     int update(int a, int b, int node = 1, int l = 0, int r = n-1) {
199         if (b < l || r < a) return tree[node];
02b         if (a <= l && r <= b) {
fd1             apply_flip(node, l, r);
14d             return tree[node];
694         }
020         prop(node, l, r);
ee4         int m = (l + r)/2;
2e4         return tree[node] = update(a, b, 2*node, l, m) + update(a, b, 2*node+1, m+1, r);
4dd     }
    
1fd     int range_set(int a, int b, int x, int node = 1, int l = 0, int r = n-1) {
199         if (b < l || r < a) return tree[node];
02b         if (a <= l && r <= b) {
98a             apply_set(node, l, r, x);
14d             return tree[node];
8dc         }
020         prop(node, l, r);
ee4         int m = (l + r)/2;
376         return tree[node] = range_set(a, b, x, 2*node, l, m) + range_set(a, b, x, 2*node+1, m+1, r);
ea9     }
    
22d     int point_set(int idx, int val, int node = 1, int l = 0, int r = n-1) {
893         if (l == r) {
a7d             lazy_set[node] = -1;
8f7             lazy_flip[node] = 0;
08c             if (val == -1) {
600                 valid[node] = 0;
c91                 tree[node] = 0;
3ca             } else {
a5b                 valid[node] = 1;
a24                 tree[node] = val;
d3e             }
14d             return tree[node];
718         }
020         prop(node, l, r);
ee4         int m = (l + r)/2;
dbc         if (idx <= m) point_set(idx, val, 2*node, l, m);
cac         else point_set(idx, val, 2*node+1, m+1, r);
d5f         valid[node] = valid[2*node] + valid[2*node+1];
b00         return tree[node] = tree[2*node] + tree[2*node+1];
63c     }
b96 }
\end{lstlisting}

\subsection{SegTree Iterativa}
\begin{lstlisting}
fb1 namespace SegTree {
1a8     int n;
8f8     int pot2;
562     vector<int> tree;
    
0d8     void build(int n2, int* v2) {
1e3         n = n2;
fa5         pot2 = 1;
2c6         while (pot2 < n) pot2 <<= 1;
26a         tree.assign(2 * pot2, 0);
4e2         for (int i = 0; i < n; ++i) tree[pot2 + i] = v2[i];
813         for (int i = pot2 - 1; i >= 1; --i) tree[i] = max(tree[2*i], tree[2*i+1]);
641     }
    
4ea     int query(int a, int b) {
026         if (a > b) return 0;
238         a += pot2; b += pot2;
11e         int res = 0;
cce         while (a <= b) {
16c             if (a & 1) res = max(res, tree[a++]);
410             if (!(b & 1)) res = max(res, tree[b--]);
691             a >>= 1; b >>= 1;
80e         }
b50         return res;
d03     }
    
cf6     void update(int i, int x) {
1a7         int p = i + pot2;
3bc         tree[p] = x;
8b8         p >>= 1;
f77         while (p >= 1) {
280             tree[p] = max(tree[2*p], tree[2*p+1]);
8b8             p >>= 1;
6f4         }
4ad     }
092 }
\end{lstlisting}

\subsection{SegTree Lazy}
\begin{lstlisting}
// Recursiva com Lazy Propagation
// Query: soma do range [a, b]
// Update: soma x em cada elemento do range [a, b]
//
// Complexidades:
// build - O(n)
// query - O(log(n))
// update - O(log(n))
//
// build (n, vector)
// query (int a, int b)
// update (int a, int b, int x)
// se for de multiplicacao mudar o lazy pra 1 e tirar o lazy * (r-l+1)

fb1 namespace SegTree {
d18 	int tree[4*MAX], lazy[4*MAX];
052 	int n, *v;
    
34c 	int build(int node=1, int l=0, int r=n-1) {
131 		lazy[node] = 0;
4e3 		if (l == r) return tree[node] = v[l];
ee4 		int m = (l+r)/2;
19c 		return tree[node] = build(2*node, l, m) + build(2*node+1, m+1, r);
851 	}
0d8 	void build(int n2, int* v2) {
680 		n = n2, v = v2;
6f2 		build();
acb 	}
dd7 	void prop(int node, int l, int r) {
7eb 		tree[node] += lazy[node]*(r-l+1);
17c 		if (l != r) lazy[2*node] += lazy[node], lazy[2*node+1] += lazy[node];
131 		lazy[node] = 0;
6eb 	}
a2c 	int query(int a, int b, int node=1, int l=0, int r=n-1) {
020 		prop(node, l, r);
e09 		if (a <= l and r <= b) return tree[node];
786 		if (b < l or r < a) return 0;
ee4 		int m = (l+r)/2;
1db 		return query(a, b, 2*node, l, m) + query(a, b, 2*node+1, m+1, r);
cf9 	}
c96 	int update(int a, int b, int x, int node=1, int l=0, int r=n-1) {
020 		prop(node, l, r);
9a3 		if (a <= l and r <= b) {
3bd 			lazy[node] += x;
020 			prop(node, l, r);
14d 			return tree[node];
57d 		}
ad9 		if (b < l or r < a) return tree[node];
ee4 		int m = (l+r)/2;
dcf 		return tree[node] = update(a, b, x, 2*node, l, m) + update(a, b, x, 2*node+1, m+1, r);
eca 	}
d1a }


//struct
383 struct SegTree {
d18     int tree[4*MAX], lazy[4*MAX];
1a8     int n;
    
037     void build(int node, int l, int r) {
131         lazy[node] = 0;
c91         tree[node] = 0;
8ce         if (l == r) return;
ee4         int m = (l+r)/2;
d58         build(2*node, l, m);
537         build(2*node+1, m+1, r);
bb1     }
    
5c7     void init(int _n) {
8fd         n = _n;
224         build(1, 0, n-1);
2e4     }
    
dd7     void prop(int node, int l, int r) {
996         if (lazy[node]) {
3d0             tree[node] = lazy[node];
579             if (l != r) {
479                 lazy[2*node] = lazy[node];
8d6                 lazy[2*node+1] = lazy[node];
98e             }
131             lazy[node] = 0;
9c8         }
f1d     }
    
7b9     void update(int a, int b, int x, int node=1, int l=0, int r=MAX-1) {
967         if (node == 1) r = n-1;
020         prop(node, l, r);
656         if (a > r || b < l) return;
02b         if (a <= l && r <= b) {
12e             lazy[node] = x;
020             prop(node, l, r);
505             return;
34c         }
ee4         int m = (l+r)/2;
c7a         update(a, b, x, 2*node, l, m);
d1f         update(a, b, x, 2*node+1, m+1, r);
7cf         tree[node] = max(tree[2*node], tree[2*node+1]);
20d     }
    
2e7     int query(int a, int b, int node=1, int l=0, int r=MAX-1) {
967         if (node == 1) r = n-1;
020         prop(node, l, r);
d9e         if (a > r || b < l) return 0;
675         if (a <= l && r <= b) return tree[node];
ee4         int m = (l+r)/2;
1e7         return max(query(a, b, 2*node, l, m), query(a, b, 2*node+1, m+1, r));
c35     }
5df };
\end{lstlisting}

\subsection{SegTree Update}
\begin{lstlisting}
// Recursiva sem Lazy Propagation
// Query: soma do range [a, b]
// Update: soma x em cada elemento do range [a, b]
//
fb1 namespace SegTree {
34a     int tree[4*MAX];
052     int n, *v;
    
34c     int build(int node=1, int l=0, int r=n-1) {
4e3         if (l == r) return tree[node] = v[l];
ee4         int m = (l+r)/2;
19c         return tree[node] = build(2*node, l, m) + build(2*node+1, m+1, r);
913     }
    
0d8     void build(int n2, int* v2) {
680         n = n2, v = v2;
6f2         build();
acb     }
       
a2c     int query(int a, int b, int node=1, int l=0, int r=n-1) {
786         if (b < l or r < a) return 0;
e09         if (a <= l and r <= b) return tree[node];
ee4         int m = (l+r)/2;
1db         return query(a, b, 2*node, l, m) + query(a, b, 2*node+1, m+1, r);
09b     }
        
eb3     void update(int i, int x, int node=1, int l=0, int r=n-1) {
893         if (l == r) {
d6d             tree[node] = x;
505             return;
b74         }
ee4         int m = (l+r)/2;
3df         if (i <= m) update(i, x, 2*node, l, m);
0a9         else update(i, x, 2*node+1, m+1, r);
8fb         tree[node] = tree[2*node] + tree[2*node+1];
dad     }
ce6 }
\end{lstlisting}

\subsection{Sparse Table}
\begin{lstlisting}
// Resolve RMQ
//
// Complexidades:
// build - O(n log(n))
// query - O(1)

cca namespace sparse {
d28 	int m[LOG+1][MAX], n;
61c 	void build(int n2, int* v) {
1e3 		n = n2;
78e 		for (int i = 0; i < n; i++) m[0][i] = v[i];
a1c 		for (int j = 1; (1<<j) <= n; j++) for (int i = 0; i+(1<<j) <= n; i++)
5d5 			m[j][i] = min(m[j-1][i], m[j-1][i+(1<<(j-1))]);
cae 	}
4ea 	int query(int a, int b) {
ee5 		int j = __builtin_clz(1) - __builtin_clz(b-a+1);
dc3 		return min(m[j][a], m[j][b-(1<<j)+1]);
fba 	}
3e1 }

67a template <typename T>
7e9 struct SparseTable {
5a6 	vector<vector<T>> m;
    	
407 	SparseTable(vector<T>& v) {
3d2 		int n = v.size();
a12 		int max_log = 32 - __builtin_clz(n);
d47 		m.assign(max_log, vector<T>(n));
78e 		for (int i = 0; i < n; i++) m[0][i] = v[i];
0de 		for (int j = 1; j < max_log; j++) {
ed9 			for (int i = 0; i + (1<<j) <= n; i++) {
5d5 				m[j][i] = min(m[j-1][i], m[j-1][i+(1<<(j-1))]);
642 			}
a87 		}
319 	}
    	
0ad 	T query(int a, int b) {
ee5 		int j = __builtin_clz(1) - __builtin_clz(b-a+1);
dc3 		return min(m[j][a], m[j][b-(1<<j)+1]);
c89 	}
4a5 };

// Sparse Table para Bitmasks
// Utiliza OR bit a bit (|) para combinar conjuntos ou atributos de um intervalo.
// Muito util para problemas com restricoes pequenas (ex: K <= 64), permitindo
// substituir K Sparse Tables por apenas 1 (ex: usando unsigned long long).
67a template <typename T>
0ff struct SparseTableMask {
5a6 	vector<vector<T>> m;
    	
b3e 	SparseTableMask(vector<T>& v) {
3d2 		int n = v.size();
a12 		int max_log = 32 - __builtin_clz(n);
d47 		m.assign(max_log, vector<T>(n));
78e 		for (int i = 0; i < n; i++) m[0][i] = v[i];
0de 		for (int j = 1; j < max_log; j++) {
ed9 			for (int i = 0; i + (1<<j) <= n; i++) {
f66 				m[j][i] = (m[j-1][i] | m[j-1][i+(1<<(j-1))]);
776 			}
e71 		}
7f1 	}
    	
0ad 	T query(int a, int b) {
ee5 		int j = __builtin_clz(1) - __builtin_clz(b-a+1);
8df 		return (m[j][a] | m[j][b-(1<<j)+1]);
7ed 	}
991 };
\end{lstlisting}



%%%%%%%%%%%%%%%%%%%%
%
% Primitivas
%
%%%%%%%%%%%%%%%%%%%%

\section{Primitivas}

\subsection{Aritmética Modular}
\begin{lstlisting}
429 template<int p> struct mod_int {
007 	int expo(int b, int e) {
856 		int ret = 1;
c87 		while (e) {
cad 			if (e % 2) ret = ret * b % p;
9d2 			e /= 2, b = b * b % p;
c42 		}
edf 		return ret;
170 	}
572 	int inv(int b) { return expo(b, p-2); }
     
4d7 	using m = mod_int;
d93 	int v;
fe0 	mod_int() : v(0) {}
cde 	mod_int(int v_) {
019 		if (v_ >= p or v_ <= -p) v_ %= p;
bc6 		if (v_ < 0) v_ += p;
2e7 		v = v_;
e1d 	}
74d 	m& operator +=(const m& a) {
2fd 		v += a.v;
ba5 		if (v >= p) v -= p;
357 		return *this;
c8b 	}
eff 	m& operator -=(const m& a) {
8b4 		v -= a.v;
cc8 		if (v < 0) v += p;
357 		return *this;
f8d 	}
4c4 	m& operator *=(const m& a) {
830 		v = v * int(a.v) % p;//tirar o int quando for usar define int long long
357 		return *this;
623 	}
3f9 	m& operator /=(const m& a) {
5d6 		v = v * inv(a.v) % p;
357 		return *this;
62d 	}
d65 	m operator -(){ return m(-v); }
e8f 	m& operator ^=(int e) {
06d 		if (e < 0) {
6e2 			v = inv(v);
00c 			e = -e;
275 		}
284 		v = expo(v, e);
    		// possivel otimizacao:
    		// cuidado com 0^0
    		// v = expo(v, e%(p-1)); 
357 		return *this;
b6d 	}
423 	bool operator ==(const m& a) { return v == a.v; }
69f 	bool operator !=(const m& a) { return v != a.v; }
     
1c6 	friend istream& operator >>(istream& in, m& a) {
9bf 		int val; in >> val;
d48 		a = m(val);
091 		return in;
a6f 	}
44f 	friend ostream& operator <<(ostream& out, m a) {
5a0 		return out << a.v;
214 	}
399 	friend m operator +(m a, m b) { return a += b; }
f9e 	friend m operator -(m a, m b) { return a -= b; }
9c1 	friend m operator *(m a, m b) { return a *= b; }
51b 	friend m operator /(m a, m b) { return a /= b; }
937 	friend m operator ^(m a, int e) { return a ^= e; }
1a1 };
 
b49 typedef mod_int<(int)MOD> mint;

e9d int modiv(int a, int b){
659     return (((a % MOD )* (fastExpo(b, MOD-2) % MOD)) % MOD);
e50 }
\end{lstlisting}

\subsection{Geometria Double}
\begin{lstlisting}
c83 typedef double ld;
e3b const ld DINF = 1e18;
43a const ld pi = acos(-1.0);
107 const ld eps = 1e-9;

b32 #define sq(x) ((x)*(x))

d97 bool eq(ld a, ld b) {
fe4 	return fabs(a - b) <= eps;
dac }

be5 struct point { // ponto
c1e 	ld x, y;
018 	point(ld x_ = 0, ld y_ = 0) : x(x_), y(y_) {}
be2 	bool operator < (const point p) const {
059 		if (!eq(x, p.x)) return x < p.x;
f98 		if (!eq(y, p.y)) return y < p.y;
bb3 		return 0;
ab8 	}
418 	bool operator == (const point p) const {
ed0 		return eq(x, p.x) and eq(y, p.y);
4a1 	}
460 	point operator + (const point p) const { return point(x+p.x, y+p.y); }
a2f 	point operator - (const point p) const { return point(x-p.x, y-p.y); }
807 	point operator * (const ld c) const { return point(x*c  , y*c  ); }
669 	point operator / (const ld c) const { return point(x/c  , y/c  ); }
986 	ld operator * (const point p) const { return x*p.x + y*p.y; }//dot
d94 	ld operator ^ (const point p) const { return x*p.y - y*p.x; }//cross
b2b 	friend istream& operator >> (istream& in, point& p) {
e37 		return in >> p.x >> p.y;
d70 	}
f70 	friend ostream& operator << (ostream& os, const point& p) {
d80 		return os << "(" << p.x << ", " << p.y << ")";
322 	}
e09 };

b3a struct line { // reta
b1f 	point p, q;
0d6 	line() {}
50a 	line(point p_, point q_) : p(p_), q(q_) {}
8d7 	friend istream& operator >> (istream& in, line& r) {
4cb 		return in >> r.p >> r.q;
858 	}
ebb };

ee4 ld dist(point p, point q) { // distancia
5f3 	return hypot(p.y - q.y, p.x - q.x);
23c }

06b ld dist2(point p, point q) { // quadrado da distancia
f24 	return sq(p.x - q.x) + sq(p.y - q.y);
2ef }

f5c ld norm(point v) { // norma do vetor
923 	return dist(point(0, 0), v);
568 }

0cf ld angle(point v) { // angulo do vetor com o eixo x
587 	ld ang = atan2(v.y, v.x);
6f8 	if (ang < 0) ang += 2*pi;
19c 	return ang;
8c4 }

3ca ld sarea(point p, point q, point r) { // area com sinal
606 	return ((q-p)^(r-q))/2;
5b6 }

6df bool col(point p, point q, point r) { // se p, q e r sao colin.
e7d 	return eq(sarea(p, q, r), 0);
20f }

c65 bool ccw(point p, point q, point r) { // se p, q, r sao ccw
fa7 	return sarea(p, q, r) > eps;
20e }

4cf point rotate(point p, ld th) { // rotaciona o ponto th radianos
4f9 	return point(p.x * cos(th) - p.y * sin(th), p.x * sin(th) + p.y * cos(th));
828 }

5cf point rotate90(point p) { // rotaciona 90 graus
128 	return point(-p.y, p.x);
c3b }

// retorna t tal que t*v pertence a reta r
849 ld get_t(point a, line r) {
ac5 	return (r.p^r.q) / ((r.p-r.q)^a);
0cf }

edc bool isvert(line r) { // se r eh vertical
87d 	return eq(r.p.x, r.q.x);
0fb }

9d8 point proj(point p, line r) { // projecao do ponto p na reta r
bea 	if (r.p == r.q) return r.p;
97a 	r.q = r.q - r.p; p = p - r.p;
ead 	point proj = r.q * ((p*r.q) / (r.q*r.q));
2cd 	return proj + r.p;
d37 }

5d2 bool isinseg(point p, line r) { // se p pertence ao seg de r
517     point a = r.p - p, b = r.q - p;
b04     return eq((a ^ b), 0) and (a * b) < eps;
da8 }

956 point inter(line r, line s) { // r inter s
b99 	if (eq((r.p - r.q) ^ (s.p - s.q), 0)) return point(DINF, DINF);
205 	r.q = r.q - r.p, s.p = s.p - r.p, s.q = s.q - r.p;
543 	return r.q * get_t(r.q, s) + r.p;
6c0 }

676 bool interseg(line r, line s) { // se o seg de r intersecta o seg de s
2b6 	if (isinseg(r.p, s) or isinseg(r.q, s) or isinseg(s.p, r) or isinseg(s.q, r)) return 1;
085 	return ccw(r.p, r.q, s.p) != ccw(r.p, r.q, s.q) and ccw(s.p, s.q, r.p) != ccw(s.p, s.q, r.q);
359 }

9f2 ld linepoint(point p, line r) { // distancia do ponto a reta
89a 	return 2 * abs(sarea(p, r.p, r.q)) / dist(r.p, r.q);
389 }

639 ld segpoint(line r, point p) { // distancia do ponto ao seg
73d 	if ((r.q - r.p)*(p - r.p) < 0) return dist(r.p, p);
951 	if ((r.p - r.q)*(p - r.q) < 0) return dist(r.q, p);
199 	return linepoint(p, r);
e86 }

ccc ld segseg(point a, point b, point c, point d) {
102 	line ab(a,b);
001 	line cd(c,d);
05c     if (interseg(ab, cd)) return 0;
494 	return min({segpoint(ab,c), segpoint(ab,d), segpoint(cd,a), segpoint(cd,b)});
3d6 }

//angulo interno entre duas linhas, mid e o ponto de intersecao entre elas
b29 ld inner_angle(point p, point q, point mid){
ac0     point p1 = (p-mid);
325     point p2 = (q-mid);
7a8     ld rad = atan2(abs(p1 ^ p2), p1 * p2);//angulo em rad
370     ld deg = rad* 180.0 / M_PI;
816     return deg;
7ec }

//linha que corta o segmento AB no meio perpendicularmente
a8a line mediatriz(point p, point q) {
da4     point mid = (p + q) / 2.0;
b5e     point vec = q - p;
5e7     point p_vec = rotate90(vec);
b03     return line(mid, mid + p_vec);
5b8 }

//bissetriz interna e externa
//bissetriz = linha que forma angulos iguais com outras duas linhas
28b pair<line, line> bissetriz(line r, line s) {
255     if (eq((r.p - r.q) ^ (s.p - s.q), 0)) return {};
ba5 	point i = inter(r,s);
1d9 	point v1 = r.p - i;
4c0     point v2 = s.p - i;
27e 	if (norm(v1) < eps) v1 = r.q - i;
672     if (norm(v2) < eps) v2 = s.q - i;
44a     v1 = v1 / norm(v1); v2 = v2 / norm(v2);
675     line interna = line(i, i + (v1 + v2));
8a6     line externa = line(i, i + (v1 - v2));
7fa     return {interna, externa};
d22 }

a0a point findcircle(point p, point q, point r){//acha o centro do circulo que encosta nos 3 pontos
55f 	if (col(p,q,r)) return point(DINF,DINF);
9b6 	line mid1 = mediatriz(p,q);
f9a 	line mid2 = mediatriz(q,r);
294 	return inter(mid1, mid2);
d41 }

a39 ld polarea(vector<point> v){ // area do poligono
9c5 	ld ret = 0;
f14 	for (int i = 0; i < v.size(); i++){
d0f 		ret += sarea(point(0, 0), v[i], v[(i + 1) % v.size()]);
8cb 	}
d03 	return abs(ret);
22a }

// se o ponto ta dentro do poligono: retorna 0 se ta fora,
// 1 se ta no interior e 2 se ta na borda
fb6 int inpol(vector<point>& v, point p) { // O(n)
8de 	int qt = 0;
f14 	for (int i = 0; i < v.size(); i++) {
bda 		if (p == v[i]) return 2;
6af 		int j = (i+1)%v.size();
e38 		if (eq(p.y, v[i].y) and eq(p.y, v[j].y)) {
97f 			if ((v[i]-p)*(v[j]-p) < eps) return 2;
5e2 			continue;
48b 		}
388 		bool baixo = v[i].y+eps < p.y;
464 		if (baixo == (v[j].y+eps < p.y)) continue;
366 		auto t = (p-v[i])^(v[j]-v[i]);
1b4 		if (eq(t, 0)) return 2;
839 		if (baixo == (t > eps)) qt += baixo ? 1 : -1;
d13 	}
b84 	return qt != 0;
1c7 }

603 ld angle(const point &p, const point &ref) { // angulo do ponto p em relacao ao ponto ref
f5b     point v = p - ref;
587     ld ang = atan2(v.y, v.x);
6f8     if (ang < 0) ang += 2*pi;
19c     return ang;
220 }

//ordena os pontos de acordo com o angulo
34e void polarSort(vector<point> &v, point &ref) {
0f8     sort(v.begin(), v.end(), [&ref](point a, point b) {
9dd         return angle(a, ref) < angle(b, ref);
17a     });
060 }

971 vector<point> convex_hull(vector<point> v) { // convex hull - O(n log(n))
fca 	sort(v.begin(), v.end());
d76 	v.erase(unique(v.begin(), v.end()), v.end());
52d 	if (v.size() <= 1) return v;
ce1 	vector<point> l, u;
f14 	for (int i = 0; i < v.size(); i++) {
fb2 		while (l.size() > 1 and !ccw(l.end()[-2], l.end()[-1], v[i]))
364 			l.pop_back();
c35 		l.push_back(v[i]);
58e 	}
3e9 	for (int i = v.size() - 1; i >= 0; i--) {
f19 		while (u.size() > 1 and !ccw(u.end()[-2], u.end()[-1], v[i]))
7a8 			u.pop_back();
a95 		u.push_back(v[i]);
0b8 	}
cfc 	l.pop_back(); u.pop_back();
2e4 	for (point i : u) l.push_back(i);
792 	return l;
9e5 }

//comprimento do arco dado o raio e o angulo
//double comprimento = (graus * M_PI * raio) / 180.0;
//double comprimento = rad * raio;
\end{lstlisting}

\subsection{Geometria Int}
\begin{lstlisting}
b2a struct pt {
e91     int x, y;
df1     pt(int x_ = 0, int y_ = 0) : x(x_), y(y_) {}
    
5bc     bool operator < (const pt p) const {
95a 		if (x != p.x) return x < p.x;
89c 		return y < p.y;
dcd 	}
    
a83 	bool operator == (const pt p) const {
d74 		return x == p.x and y == p.y;
7b4 	}
    
7af     pt operator + (const pt& p) const { return pt(x + p.x, y + p.y); }
b23     pt operator - (const pt& p) const { return pt(x - p.x, y - p.y); }
477     template<typename U> 
60f     pt operator * (const U c) const { return pt(x * c, y * c); }
abb     int operator * (const pt& p) const { return x * p.x + y * p.y; }
476     int operator ^ (const pt& p) const { return x * p.y - y * p.x; }
    
e45     friend istream& operator >> (istream& in, pt& p) { return in >> p.x >> p.y; }
03c     friend ostream& operator << (ostream& os, const pt& p) { return os << "(" << p.x << ", " << p.y << ")"; }
b4a };

b3a struct line {
730 	pt p, q;
0d6 	line() {}
4b8 	line(pt p_, pt q_) : p(p_), q(q_) {}
    
8d7     friend istream& operator >> (istream& in, line& r) {
4cb 		return in >> r.p >> r.q;
858 	}
    
6b7     friend ostream& operator << (ostream& os, line& r) {
e8e 		return os << "[" << r.p << " -> " << r.q << "]";
bbb 	}
3b6 };

46d int sq(pt p){
8a7     return p.x*p.x + p.y*p.y;
b0e }

//checa se dois vetores sao perpendiculares
6b8 bool perpendicular(pt v, pt w){
350     return (v*w == 0);
99e }

//ccw, 0 colineares, -1 r a direita, 1 r a esquerda
f2b int orientation(pt p, pt q, pt r){
c5b     int val = ((q-p) ^ (r-p));
db6     if (!val) return 0;
326     return (val > 0) ? 1 : -1;
dd6 }

//checa se um ponto esta contido no angulo do segmento pq e pr
e82 bool onAngle (pt p, pt q, pt r, pt g){
999     if (!orientation(p,q,r)) return false;
673     if (orientation(p,q,r) > 0) swap(q,r);
d46     return orientation(p,q,g) >= 0 && orientation(p,r,g) <= 0;
edb }

// quadrante de um ponto
c31 int quad(pt p) {
dbb 	return (p.x<0)^3*(p.y<0);
fcf }

//dist quadrada de dois pontos
9ea int dist2(pt p, pt q){
fda     return sq(p - q);
fb2 }

// 2 * area com sinal
a15 int sarea2(pt p, pt q, pt r) {
586 	return (q-p)^(r-q);
9ef }

e32 bool col(pt p, pt q, pt r) { // se p, q e r sao colin.
034 	return sarea2(p, q, r) == 0;
a08 }

// 2 * area do poligono
fc1 int polarea2(vector<pt> v) {
7c9 	int ret = 0;
c6e 	for (int i = 0; i < v.size(); i++)
532 		ret += sarea2(pt(0, 0), v[i], v[(i + 1) % v.size()]);
d03 	return abs(ret);
87e }

//checa se o poligono e convexo
f12 bool isConvex(vector<pt> p) {
ff7     bool hasPos = false, hasNeg = false;
ed0     for (int i = 0, n = p.size(); i < n; i++) {
568         int o = orientation(p[i], p[(i+1)%n], p[(i+2)%n]);
f42         if (o > 0) hasPos = true;
162         if (o < 0) hasNeg = true;
455     }
f61     return !(hasPos && hasNeg);
61d }

ae2 bool half(pt p) { // true if in blue half
184     assert(p.x != 0 || p.y != 0); // the argument of (0,0) isundefined
a12     return p.y > 0 || (p.y == 0 && p.x < 0);
a70 }

//ordena os pontos de acordo com o angulo
61f void polarSort(vector<pt> &v) {
2f2     sort(v.begin(), v.end(), [](pt v, pt w) {
8ba         return make_tuple(half(v), 0, sq(v)) <
99d         make_tuple(half(w), (v^w), sq(w));
486     });
34d }

//rotaciona o vetor 90 graus
ab1 pt rotate90(pt p) {
a0d 	return pt(-p.y, p.x);
e4a }

//retorna uma linha perpendicular a linha s que passa pelo ponto p
9c9 line perpthrough(line s, pt p){
2d4     return {p, p + rotate90(s.p - s.q)};
299 }

// se p esta na linha de r
985 bool onLine(line r, pt p){
489     return col(r.p, r.q, p);
        // return ((r.q - r.p) ^ (p - r.p)) == 0;
ea4 }

// se p pertence ao seg de r
63e bool isinseg (line r, pt p) {
eb7 	pt a = r.p - p;
e99     pt b = r.q - p;
2ac 	return (a ^ b) == 0 and (a * b) <= 0;
336 }

147 bool intersect(line r, line s){
0d5     int o1 = orientation(r.p, r.q, s.p);
f92     int o2 = orientation(r.p, r.q, s.q);
033     int o3 = orientation(s.p, s.q, r.p);
885     int o4 = orientation(s.p, s.q, r.q);
    
bb3     if (o1 != o2 && o3 != o4) return true;
    
171     if (o1 == 0 && isinseg(r, s.p)) return true;
0f8     if (o2 == 0 && isinseg(r, s.q)) return true;
082     if (o3 == 0 && isinseg(s, r.p)) return true;
718     if (o4 == 0 && isinseg(s, r.q)) return true;
    
d1f     return false;
e86 }

// numero de pontos inteiros no segmento
fae int segpts(line r) {
508 	return 1 + gcd(abs(r.p.x - r.q.x), abs(r.p.y - r.q.y));
28a }

// se o ponto ta dentro do poligono: retorna 0 se ta fora,
// 1 se ta no interior e 2 se ta na borda
8e7 int inpol(vector<pt>& v, pt p){ // O(n)
8de 	int qt = 0;
f14 	for (int i = 0; i < v.size(); i++) {
bda 		if (p == v[i]) return 2;
6af 		int j = (i+1)%v.size();
cc6 		if (p.y == v[i].y and p.y == v[j].y) {
547 			if ((v[i]-p)*(v[j]-p) <= 0) return 2;
5e2 			continue;
b47 		}
78c 		bool baixo = v[i].y < p.y;
057 		if (baixo == (v[j].y < p.y)) continue;
366 		auto t = (p-v[i])^(v[j]-v[i]);
2ad 		if (!t) return 2;
0bb 		if (baixo == (t > 0)) qt += baixo ? 1 : -1;
9cf 	}
b84 	return qt != 0;
afd }
\end{lstlisting}

\subsection{nclude <bits/stdc++.h>}
\begin{lstlisting}

613 #define all(x) x.begin(), x.end()
1c9 #define rall(x) x.rbegin(), x.rend()
42a #define endl '\n'
a7f #define int long long
d40 #define ld long double

7f0 namespace dbg {
f5b     constexpr const char* RESET      = "\033[0m";
c6b     constexpr const char* BOLD_BLUE  = "\033[1;34m";
2d5     constexpr const char* YELLOW     = "\033[33m";
c5b     constexpr const char* BOLD_WHITE = "\033[1;37m";
    
1cc     template<typename T, typename = void> struct is_container : false_type {};
d90     template<typename T> struct is_container<T, void_t<typename T::value_type>> : bool_constant<!is_same_v<T, string> && !is_same_v<T, string_view>> {};
    
9df     inline ostream& print_atom(ostream& os, bool b)         { return os << (b ? "true" : "false"); }
de4     inline ostream& print_atom(ostream& os, char c)         { return os << '\'' << c << '\''; }
c4a     inline ostream& print_atom(ostream& os, const string& s){ return os << '"' << s << '"'; }
02d     inline ostream& print_atom(ostream& os, string_view s)  { return os << '"' << s << '"'; }
9aa     inline ostream& print_atom(ostream& os, const char* s)  { return os << '"' << s << '"'; }
    
8c4     template<typename T, enable_if_t<!is_container<T>::value, int> = 0>
b90     ostream& print_atom(ostream& os, const T& x) { return os << x; }
    
620     template<typename T1, typename T2>
853     ostream& print_atom(ostream& os, const pair<T1, T2>& p) {
881         os << '{'; print_atom(os, p.first); os << ", "; print_atom(os, p.second); return os << '}';
3f2     }
6d2     template<typename... T>
c49     ostream& print_atom(ostream& os, const tuple<T...>& t) {
b6a         os << '{';
96d         apply([&os](auto const&... args) {
890             size_t n = 0;
f00             ((print_atom(os, args), os << (++n != sizeof...(T) ? ", " : "")), ...);
8ee         }, t);
925         return os << '}';
32f     }
    
f91     template<typename T, typename C>
127     ostream& print_atom(ostream& os, stack<T, C> s) {
ba3         os << '{'; bool f = true; while (!s.empty()) { if (!f) os << ", "; print_atom(os, s.top()); s.pop(); f = false; } return os << '}';
1fc     }
f91     template<typename T, typename C>
732     ostream& print_atom(ostream& os, queue<T, C> q) {
b98         os << '{'; bool f = true; while (!q.empty()) { if (!f) os << ", "; print_atom(os, q.front()); q.pop(); f = false; } return os << '}';
c48     }
bab     template<typename T, typename C, typename Cmp>
12d     ostream& print_atom(ostream& os, priority_queue<T, C, Cmp> q) {
7f3         os << '{'; bool f = true; while (!q.empty()) { if (!f) os << ", "; print_atom(os, q.top()); q.pop(); f = false; } return os << '}';
aa9     }
    
383     template<typename C, enable_if_t<is_container<C>::value, int> = 0>
647     ostream& print_atom(ostream& os, const C& v) {
b64         os << '{'; bool f = true;
3c0         for (const auto& x : v) { if (!f) os << ", "; print_atom(os, x); f = false; }
b9e         return os << '}';
fa4     }
    
383     template<typename C, enable_if_t<is_container<C>::value, int> = 0>
132     ostream& operator<<(ostream& os, const C& v)            { return print_atom(os, v); }
620     template<typename T1, typename T2>
51b     ostream& operator<<(ostream& os, const pair<T1, T2>& p) { return print_atom(os, p); }
6d2     template<typename... T>
86c     ostream& operator<<(ostream& os, const tuple<T...>& t)  { return print_atom(os, t); }
    
eb5     inline size_t split_arg(string_view s) {
cea         int depth = 0;
9f5         for (size_t i = 0; i < s.size(); ++i) {
827             char c = s[i];
47c             if (c == '"' || c == '\'') {
d67                 char q = c; ++i;
19f                 while (i < s.size() && s[i] != q) { if (s[i] == '\\' && i + 1 < s.size()) ++i; ++i; }
2c2             } else if (c == '(' || c == '[' || c == '{' || c == '<') ++depth;
edb             else if (c == ')' || c == ']' || c == '}' || c == '>') --depth;
9c3             else if (c == ',' && depth == 0) return i;
564         }
ddb         return string_view::npos;
cc4     }
    
720     inline void debug_out(string_view) { cerr << '\n' << flush; }
2dc     template<typename H, typename... T>
943     void debug_out(string_view s, const H& h, const T&... t) {
940         size_t cpos = split_arg(s);
77a         auto name = s.substr(0, cpos);
df4         while (!name.empty() && (name.front() == ' ' || name.front() == '\t')) name.remove_prefix(1);
3d4         while (!name.empty() && (name.back()  == ' ' || name.back()  == '\t')) name.remove_suffix(1);
0d8         cerr << YELLOW << name << RESET << " = " << BOLD_WHITE;
5c7         print_atom(cerr, h) << RESET;
c4a         if constexpr (sizeof...(t) > 0) {
57b             cerr << ", ";
036             debug_out(cpos == string_view::npos ? string_view{} : s.substr(cpos + 1), t...);
fad         } 
419         else cerr << '\n' << flush;
1d6     }
bec }
e4e using namespace dbg;

344 #define DEBUG

733 #if defined(DEBUG)
769     #define winton (void)0
240     #define debug(...) do { cerr << BOLD_BLUE << "[" << __func__ << ":" << __LINE__ << "] " << RESET; dbg::debug_out(#__VA_ARGS__, __VA_ARGS__); } while (0)
8c1 #else
62e     #define winton ios_base::sync_with_stdio(false),cin.tie(NULL),cout.tie(NULL)
3ae     #define debug(...) ((void)0)
f2e #endif

b2a struct pt {
e91     int x, y;
df1     pt(int x_ = 0, int y_ = 0) : x(x_), y(y_) {}
        
5bc     bool operator < (const pt p) const {
95a         if (x != p.x) return x < p.x;
89c 		return y < p.y;
dcd 	}
        
a83 	bool operator == (const pt p) const {
d74         return x == p.x and y == p.y;
7b4 	}
        
7af     pt operator + (const pt& p) const { return pt(x + p.x, y + p.y); }
b23     pt operator - (const pt& p) const { return pt(x - p.x, y - p.y); }
477     template<typename U>
60f     pt operator * (const U c) const { return pt(x * c, y * c); }
818     auto operator * (const pt& p) const { return (int)x * p.x + (int)y * p.y; }
9a6     auto operator ^ (const pt& p) const { return (int)x * p.y - (int)y * p.x; }
        
e45     friend istream& operator >> (istream& in, pt& p) { return in >> p.x >> p.y; }
03c     friend ostream& operator << (ostream& os, const pt& p) { return os << "(" << p.x << ", " << p.y << ")"; }
        
df0 };


0dd signed main(){
5d6     winton;
6cb     pt a, b, c;;
e22     cin >> a >> b >> c;
56d     pt p(2,3);
00f     pt q(4,1);
    
129     debug(a-c, b-c);
57c     int area = (a-c) ^ (b-c);
f4b     int area2 = (p) ^ (q);
f31     debug(area2, area);
    
2dc }
\end{lstlisting}

\subsection{ruct matrix : vector<vector<int>> {}
\begin{lstlisting}

df8     matrix(int n_, int m_, bool ident = false) : vector<vector<int>>(n_, vector<int>(m_, 0)), n(n_), m(m_) {
94e         if (ident) {
df7             assert(n == m);
a89 			for (int i = 0; i < n; i++) (*this)[i][i] = 1;
359 		}
456 	}

483     matrix(const vector<vector<int>>& c) : vector<vector<int>>(c),
a3d     n(c.size()), m(c[0].size()) {}

d01     matrix operator*(const matrix &a) const {
eda         assert(m == a.n);
ca0         matrix res(n, a.m);
603         for (int i = 0; i < n; i++){
cf4             for (int j = 0; j < a.m; j++){
917                 for (int k = 0; k < m; k++){
956                     res[i][j] = (res[i][j] + ((*this)[i][k] * a[k][j]) % MOD) % MOD;
368                 }
830             }
a00         }
b50         return res;
3b4     }

3d2     matrix operator^(int e) const {
cd2         matrix res(n,n,true);
2de         matrix base = *this;
c87         while (e) {
d4a             if (e&1) res = res * base;
6c1             base = base * base;
cc2             e >>= 1;
ece         }
b50         return res;
81d     }
214 };
\end{lstlisting}



%%%%%%%%%%%%%%%%%%%%
%
% Problemas
%
%%%%%%%%%%%%%%%%%%%%

\section{Problemas}

\subsection{}
\begin{lstlisting}
41d inversion pairs before each removal. The number of inversion pairs of an array A is the number of
a9e ordered pairs (i, j) such that i < j and A[i] > A[j].
c4c */

2b7 #include <bits/stdc++.h>
ca4 using namespace std;

613 #define all(x) x.begin(), x.end()
1c9 #define rall(x) x.rbegin(), x.rend()
42a #define endl '\n'
a7f #define int long long
d40 #define ld long double

7f0 namespace dbg {
404     const char* const RESET     = "\033[0m";
646     const char* const BOLD_BLUE = "\033[1;34m";
d39     const char* const YELLOW    = "\033[33m";
e1d     const char* const BOLD_WHITE= "\033[1;37m";
    
620     template<typename T1, typename T2>
c1f     ostream& operator<<(ostream& os, const pair<T1, T2>& p);
    
876     template<typename T_container, typename T = typename enable_if<!is_same_v<T_container, string> && !is_same_v<T_container, string_view>, typename T_container::value_type>::type>
dab     ostream& operator<<(ostream& os, const T_container& v) {
b6a         os << '{';
ce4         bool first = true;
dbd         for (const T& x : v) { os << (first ? "" : ", ") << x, first = false; }
ff7         return os << '}';
1d7     }
    
620     template<typename T1, typename T2>
598     ostream& operator<<(ostream& os, const pair<T1, T2>& p) { return os << '{' << p.first << ", " << p.second << '}'; }
    
3a4     void debug_out(string_view) { cerr << endl; }
2dc     template<typename H, typename... T>
dc3     void debug_out(string_view s, H h, T... t) {
196         auto cpos = s.find(',');
53e         cerr << YELLOW << s.substr(0, cpos) << RESET << " = ";
891         cerr << BOLD_WHITE << h << RESET;
c4a         if constexpr (sizeof...(t) > 0) {
57b             cerr << ", ";
1cb             auto nx = s.find_first_not_of(" \t\n\r", cpos + 1);
937             debug_out(s.substr(nx), t...);
7eb         } else {
016             cerr << endl;
bf4         }
6ed     }
b19 } 
e4e using namespace dbg;

// #define DEBUG

733 #if defined(DEBUG)
769     #define winton (void)0
c88     #define debug(...) cerr << BOLD_BLUE << "[" << __func__ << ":" << __LINE__ << "]" << RESET << " "; debug_out(#__VA_ARGS__, __VA_ARGS__)
8c1 #else
62e     #define winton ios_base::sync_with_stdio(false),cin.tie(NULL),cout.tie(NULL)
e3e     #define debug(...) (void)0
f2e #endif


a6b template<class T = int> struct bit2d {
acf 	vector<T> X;
a84 	vector<vector<T>> Y, t;
    
709 	int ub(vector<T>& v, T x) {
dde 		return upper_bound(v.begin(), v.end(), x) - v.begin();
9cc 	}
5cb 	bit2d(vector<pair<T, T>> v) {
2e1 		for (auto [x, y] : v) X.push_back(x);
fd4 		sort(X.begin(), X.end());
1ee 		X.erase(unique(X.begin(), X.end()), X.end());
    
d56 		t.resize(X.size() + 1);
d12 		Y.resize(t.size());
3d0 		sort(v.begin(), v.end(), [](auto a, auto b) {
e8f 			return a.second < b.second; });
961 		for (auto [x, y] : v) for (int i = ub(X, x); i < t.size(); i += i&-i)
b75 			if (!Y[i].size() or Y[i].back() != y) Y[i].push_back(y);
    
7c7 		for (int i = 0; i < t.size(); i++) t[i].resize(Y[i].size() + 1);
8cc 	}
    
e78 	void update(T x, T y, T v) {
2a9 		for (int i = ub(X, x); i < t.size(); i += i&-i)
cd2 			for (int j = ub(Y[i], y); j < t[i].size(); j += j&-j) t[i][j] += v;
533 	}
    
5d2 	T query(T x, T y) {
966 		T ans = 0;
c54 		for (int i = ub(X, x); i; i -= i&-i)
4fb 			for (int j = ub(Y[i], y); j; j -= j&-j) ans += t[i][j];
ba7 		return ans;
62d 	}
46d 	T query(T x1, T y1, T x2, T y2) {
fcf 		return query(x2, y2)-query(x2, y1-1)-query(x1-1, y2)+query(x1-1, y1-1);
232 	}
6a7 };

0dd signed main(){
5d6     winton;
3b9     int n, q;
83f     while(cin >> n >> q){
    
7ed         vector<int> a(n), pos(n+1);
a84         vector<pair<int,int>> plano(n);
603         for (int i = 0; i < n; i++){
788             cin >> a[i];
d1d             pos[a[i]] = i+1;
1b4             plano[i] = {i+1, a[i]};
df0         }
            
9ff         bit2d bit(plano);
2cb         for (auto [x,y] : plano) bit.update(x,y,1);
a93         int sum = 0;
e84         for (auto [x,y] : plano){
d6b             sum += bit.query(x+1, 1, n, y-1);
7db         } 
cde         debug(sum);
            
1f4         while(q--){
                // cout << "sum: " ;
d94             cout << sum << endl;
                
93d             int y;
e90             cin >> y;
874             int x = pos[y];
e04             bit.update(x,y,-1);
053             sum -= bit.query(1, y+1, x-1, n);
89a             sum -= bit.query(x+1, 1, n, y-1);
5b7         }       
7e6     }
698 }
\end{lstlisting}

\subsection{Artwork}
\begin{lstlisting}
//Grafo de varios circulos
//saber se chega do 00 ao nm sem encostar nos circulos

add int n, m, k; 
d56 struct DSU{
1a8     int n;
632     vector<int> p;
02d     vector<int> r;
a95     vector<vector<int>> encosta;
2ea     DSU(int _n, vector<vector<int>> _marcador){
8fd         n = _n;
1c4         p.resize(n);
7e4         r.resize(n);
972         encosta = _marcador;
603         for(int i = 0; i < n; i++){
f13             p[i] = i;
fd1             r[i] = 0;
606         }
68f     }
57b     int find(int x){
eab         return (p[x] == x? x : p[x] = find(p[x]));
ab9     }
440     void unite(int a, int b){
605         a = find(a), b = find(b);
d54         if(a == b) return;
dac         if(r[a] == r[b]) r[a]++;
d1b         if(r[a] > r[b]){
67e             r[a] += r[b];
264             p[b] = a;
1cd             for(int i = 0; i < 4; i++){
9e8                 encosta[a][i] |= encosta[b][i];
31b             }
e6c         }else{
2e3             r[b] += r[a];
84b             p[a] = b;
1cd             for(int i = 0; i < 4; i++){
7a5                 encosta[b][i] |= encosta[a][i];
f4c             }
701         }
15f     }
b43 };
80d int dist(int x1, int y1, int x2, int y2){
a9d     return (x1-x2)*(x1-x2) + (y1-y2)*(y1-y2);
fb0 }
0dd signed main(){
9ae     fastio;
5fe     cin >> n >> m >> k;
6ef     int pos = 1;
c56     vector<tuple<int, int, int>> v;
0bd     vector<vector<int>> marcador(k, vector<int>(4, 0));
b3d     for(int i = 0; i < k; i++){
32e         int a, b, c; cin >> a >> b >> c;
6c9         v.push_back({a, b, c});
ea0         if(a-c <= 0) marcador[i][0] = 1;
8a1         if(a+c >= n) marcador[i][1] = 1;
552         if(b-c <= 0) marcador[i][2] = 1;
dfd         if(b+c >= m) marcador[i][3] = 1;
f3a     }
    
313     DSU dsu(k, marcador);
b3d     for(int i = 0; i < k; i++){
dde         auto [x1, y1, r1] = v[i];
453         for(int j = i+1; j < k; j++){
b4a             if(j == i) continue;
f7c             auto [x2, y2, r2] = v[j];
d6c             if(dist(x1, y1, x2, y2) <= (r1 + r2)*(r1+r2)){
7fd                 dsu.unite(i, j);
823             }
88d         }
c50     }
    
b3d     for(int i = 0; i < k; i++){
a44         auto p = dsu.find(i);
00d         if(dsu.encosta[p][0] && dsu.encosta[p][1]){
792             pos = 0;
c2b             break;
873         }if(dsu.encosta[p][2] && dsu.encosta[p][3]){
792             pos = 0;
c2b             break;
fd0         }
089         if(dsu.encosta[p][0] && dsu.encosta[p][2]){
792             pos = 0;
c2b             break;
3c1         }
b6e         if(dsu.encosta[p][1] && dsu.encosta[p][3]){
792             pos = 0;
c2b             break;
aff         }
035     }
c69     if(pos) cout << "S\n";
100     else cout << "N\n";
77c }
\end{lstlisting}

\subsection{Centro de Árvore}
\begin{lstlisting}
309 vector<vector<int>> graph;
ac8 int n, dist[3][MAX];

69f void dfs(int v, int p, int id){
284     for (auto u : graph[v]){
2ab         if (u == p || dist[id][u] > 0) continue;
14a         dist[id][u] = dist[id][v] + 1;
041         dfs(u,v,id);
0ea     }
d8f }
da9 int distance(int id){
bda     int u = -1, mx = 0;
78a     for (int i = 1; i <= n; i++){
efa         if (dist[id][i] > mx){
af1             u = i;
87e             mx = dist[id][i];
da3         }
c2f     }
03f     return u;
f26 }
738 int cont(int v, int p, int d){
e40     if (d == 0) return 1;
b3d     int cur = 0;
284     for (auto u : graph[v]){
2da         if (u == p) continue;
72d         cur += cont(u,v,d-1);
075     }
75e     return cur;
2a4 }

0dd signed main () {
5d6     winton;
a68     cin >> n;
8f4     graph.resize(n+1);
6f5     for (int i = 1; i < n; i++){
ba2         int a, b;
0a8         cin >> a >> b;
8fa         graph[a].push_back(b);
4c6         graph[b].push_back(a);
080     }
c7e     dfs(1,0,0);
398     int a = distance(0);
f96     dfs(a,0,1);
924     int b = distance(1);
ee1     dfs(b,0,2);
    
ade     vector<int> center;
78a     for (int i = 1; i <= n; i++){
8b2         if (dist[1][i] + dist[2][i] == dist[1][b] && abs(dist[1][i] - dist[2][i]) <= 1){
1dc             center.push_back(i);
cef         }
7f6     }
319 }
\end{lstlisting}

\subsection{Coin Collector}
\begin{lstlisting}
4d3 vector<vector<int>> graph, rgraph;
c7b vector<set<int>> gcomp;
7b4 stack<int> st;
daa int n, com, indeg[MAX], dp[MAX], coins[MAX], visited[MAX], comp[MAX], gcoins[MAX];
 
76a void dfs(int v) {
cda 	visited[v] = 1;
284 	for (auto u : graph[v]){
e39 		if (!visited[u]) dfs(u);
6e1     }
789 	st.push(v);
850 }
 
664 void scc(int v, int c) {
cda 	visited[v] = 1;
d5f 	comp[v] = c;
af8 	for (auto u : rgraph[v]){
157 		if (!visited[u]) scc(u, c);
7bd     }
7e6 }
 
db8 void kosaraju() {
1ef 	for (int i = 0; i < n; i++) visited[i] = 0;
76b 	for (int i = 0; i < n; i++) if (!visited[i]) dfs(i);
1ef 	for (int i = 0; i < n; i++) visited[i] = 0;
3c6 	while (st.size()) {
f1a 		int u = st.top();
25a 		st.pop();
264 		if (!visited[u]) {
e25             scc(u, com++);
1b3         }
338     }
4de }
 
0dd signed main(){
2e6     int m;
aa3     cin >> n >> m;
c02     graph.resize(n);
21d     rgraph.resize(n);
     
5ed     for (int i = 0; i < n; i++) cin >> coins[i];
dd5     for (int i = 0; i < m; i++){
ba2         int a, b;
0a8         cin >> a >> b;
9c7         graph[--a].push_back(--b);
8d3         rgraph[b].push_back(a);
a20     }
75d     kosaraju();
b2b     gcomp.resize(com);
778     for(int v = 0; v < n; v++){
284         for (auto u : graph[v]){
aa8             if (comp[u] == comp[v]) continue;
            
ad6             if (gcomp[comp[v]].insert(comp[u]).second){
980                 indeg[comp[u]]++;
384             }
504         }
601     }
7c8     for(int i = 0; i < n; i++) gcoins[comp[i]] += coins[i];
425     for(int i = 0; i < n; i++) dp[i] = gcoins[i];
603     for(int i = 0; i < n; i++) {
ad3         debug(i, comp[i], gcoins[i], indeg[i]);
d2c     }
     
26a     queue<int> q;
c97     for (int i = 0; i < com; i++){
ac5         if (indeg[i] == 0) q.push(i);
21b     }
     
14d     while (!q.empty()){
b1e         int v = q.front();
833         q.pop();
1ed         for (auto u : gcomp[v]){
7bd             debug(u);
f27             dp[u] = max(dp[u], dp[v] + gcoins[u]);
e80             if (--indeg[u] == 0) q.push(u);
ee4         }
b12     }
1a4     int ans = 0;
603     for (int i = 0; i < n; i++){
1cf         ans = max(ans, dp[i]);
f56     }
886     cout << ans << endl;
187 }
\end{lstlisting}

\subsection{Creating Multiples}
\begin{lstlisting}
0dd signed main(){
5d6     winton;
b56     int b, l;
110     cin >> b >> l;
e28     vector<int> d(l);
cf7     for(auto &u : d) cin >> u;
afa     int base = b+1;
660     int resto = 0;
d00     for (int i = 0; i < l; i++) {
01a         int termo = d[i];
002         if (!((l-1-i)&1)) {
142             termo = -termo;
87e         }
5c5         resto = (resto+termo % base + base) % base;
de7     }
9ed     if (resto == 0) {
005         cout << "0 0" << endl;
bb3         return 0;
a7c     }
d00     for (int i = 0; i < l; i++) {
4eb         int reducao;
600         if ((l-1-i)&1) {
87d             reducao = resto;
0b4         } 
2b6         else reducao = base - resto;
     
aa1         if (d[i] >= reducao) {
092             cout << i + 1 << " " << d[i] - reducao << endl;
bb3             return 0;
c6f         }
8da     }
f82     cout << "-1 -1" << endl;;
960 }
\end{lstlisting}

\subsection{Divisor Analysis}
\begin{lstlisting}
// quantidade, soma e produto dos divisores

bf9 int fastExpo(int base, int exp) {
c53     int r = 1;
479     base %= MOD;
fb9     while(exp) {
727         if (exp & 1) r = r * base % MOD;
a6a         base = base * base % MOD;
ef0         exp >>= 1;
a45     }
4c1     return r;
844 }

fdd int modinverse(int n) {
562     return fastExpo(n, MOD - 2);
8ff }

e9d int modiv(int a, int b){
0e2     return (a * modinverse(b))%MOD;
656 }

0dd signed main(){
1a8     int n;
a68     cin >> n;
dd9     bool par = true;
b7a     vector<int> p(n), k(n);
f40     int cnt = 1, sum = 1, prod = 1, cnt2 = 1;
603     for (int i = 0; i < n; i++){
742         cin >> p[i] >> k[i];
9d8         if (k[i]&1) par = false;
dd3     }
603     for (int i = 0; i < n; i++){          
ed6         cnt = (cnt * (k[i]+1)) % MOD;
57a         cnt2 = (cnt2 * (k[i]+1)) % (2*(MOD-1));
            
a4a         int exp = (k[i] % (MOD-1));
db7         exp = (exp+1) % (MOD-1);
            //debug(exp);
ab0         int num = (fastExpo(p[i],exp)-1 + MOD) % MOD;
27b         int den = (p[i]-1 + MOD) % MOD;
ce9         sum = (sum * modiv(num,den)%MOD)%MOD;
a41     }
603     for (int i = 0; i < n; i++){       
b5b         if (!par){
bf2             int half = (cnt2/2)%(MOD-1);
220             int exp = (k[i]*half)%(MOD-1);
dc0             prod = (prod * (fastExpo(p[i],exp)%MOD)) % MOD;
b1f         }
4e6         else{
c85             int exp = ((k[i]/2)*cnt2)%(MOD-1);
dc0             prod = (prod * (fastExpo(p[i],exp)%MOD)) % MOD;
719         }
e4f     }
d2a     cout << cnt%MOD << ' ' << sum%MOD << ' ' << prod%MOD << endl;
a52 }
\end{lstlisting}

\subsection{Drone Control (3 Equa, 4 Vars)}
\begin{lstlisting}
// n(s) = r + s
// e(s) = (y+p-r)/2 - s
// w(s) = (y-p-r)/2 - s
// s(s) = s
 
// n(s) = a + s
// e(s) = b - s
// w(s) = c - s
// s(s) = s
 
// M = max{(|a+s|), (|b-s|), (|c-s|), (|s|)}
 
// |a+s| <= M    ->     -M-a <= s <= M-a
// |b-s| <= M    ->     -M+b <= s <= M+b
// |c-s| <= M    ->     -M+c <= s <= M+c
// |s|   <= M    ->     -M   <= s <= M
 
// achar o minimo e maximo do intervalo de s
 
// smin = max(-a,b,c,0)
// smax = max(-a,b,c,0)
// 
0dd signed main(){
45f     int q;
11b     cin >> q;
1f4     while(q--){
172         ld p, r, y;
b47         cin >> p >> r >> y;
bf3         ld a = r;
09e         ld b = (y+p-r)/2;
305         ld c = (y-p-r)/2;
     
c79         ld smin = max({-a,b,c,(ld)0.0});
34f         ld smax = min({-a,b,c,(ld)0.0});
     
e2e         ld s = (smin+smax)/2;
     
4cb         ld n = r+s;
e48         ld e = b-s;
c2d         ld w = c-s;
     
e8a         cout << fixed << setprecision(7) << n << ' ' << e << ' ' << s << ' ' << w << ' ' << endl;
bb1     }
cb3 }
\end{lstlisting}

\subsection{Hungaro}
\begin{lstlisting}
// Resolve o problema de assignment (matriz n x n)
// Colocar os valores da matriz em 'a' (pode < 0)
// assignment() retorna um par com o valor do
// assignment minimo, e a coluna escolhida por cada linha
//
// Se precisar calcular o produtorio da pra guardar o log(double(x)) no a[i][j]
// -log se for maior produtorio e +log se for o menor produtorio
// O(n^3)

a6a template<typename T> struct hungarian{
1a8     int n;
a08     vector<vector<T>> a;
f36     vector<T> u, v;
5ff     vector<int> p, way;
f1e     T inf;
    
aea     hungarian(int _n, vector<vector<T>> _a){
8fd         n = _n; 
92d         u.resize(n+1);
eac         v.resize(n+1);
181         p.resize(n+1); 
c7c         way.resize(n+1);
ad6         a = _a;
d80         inf = INT_MAX;
c00     }
    
d67    	pair<T, vector<int>> assignment() {
78a 		for (int i = 1; i <= n; i++) {
8c9 			p[0] = i;
625 			int j0 = 0;
ce7 			vector<T> minv(n+1, inf);
241 			vector<int> used(n+1, 0);
016 			do {
472 				used[j0] = true;
d24 				int i0 = p[j0], j1 = -1;
7e5 				T delta = inf;
9ac 				for (int j = 1; j <= n; j++) if (!used[j]) {
7bf 					T cur = a[i0-1][j-1] - u[i0] - v[j];
9f2 					if (cur < minv[j]) minv[j] = cur, way[j] = j0;
821 					if (minv[j] < delta) delta = minv[j], j1 = j;
4d1 				}
f63 				for (int j = 0; j <= n; j++)
2c5 					if (used[j]) u[p[j]] += delta, v[j] -= delta;
6ec 					else minv[j] -= delta;
6d4 				j0 = j1;
f4f 			} while (p[j0] != 0);
016 			do {
4c5 				int j1 = way[j0];
0d7 				p[j0] = p[j1];
6d4 				j0 = j1;
886 			} while (j0);
38d 		}
306 		vector<int> ans(n);
6db 		for (int j = 1; j <= n; j++) ans[p[j]-1] = j-1;
da3 		return make_pair(-v[0], ans);
979 	}
51f };
\end{lstlisting}

\subsection{Intervalo da String Palíndromico}
\begin{lstlisting}
//Acha o maior intervalo que pode ser reoordenado como palindrome numa string
//Pode ser adaptado para contar quantos palindromes e possivel fazer 

071 int size = (1<<26);
7cf vector<int> first(size, -1);
198 int mask = 0;
b55 first[0] = 0;
87c int ans = 1;
1cf for (int i = 0; i < s.size(); i++) {
90e     int chab = s[i] - 'a';
fe3     mask ^= (1 << chab);
208     if (first[mask] != -1) ans = max(ans, (i - first[mask])+1);
033     for (int k = 0; k < 26; ++k) {
f9c         int targetmask = mask ^ (1 << k);
dc8         if (first[targetmask] != -1) ans = max(ans, (i - first[targetmask])+1);
804     }
c26     if (first[mask] == -1) first[mask] = i+1;
b7d }
\end{lstlisting}

\subsection{Inversões}
\begin{lstlisting}
833 int countAndMerge(vector<int>& arr, int l, int m, int r) {
fd8     int n1 = m - l + 1, n2 = r - m;
e11     vector<int> left(n1), right(n2);
bac     for (int i = 0; i < n1; i++) left[i] = arr[i + l];
d74     for (int j = 0; j < n2; j++) right[j] = arr[m + 1 + j];
5e8     int res = 0, i = 0, j = 0, k = l;
40c     while (i < n1 && j < n2) {
839         if (left[i] <= right[j]) arr[k++] = left[i++];
4e6         else {
20c             arr[k++] = right[j++];
091             res += (n1 - i);
481         }
ceb     }
694     while (i < n1) arr[k++] = left[i++];
3cc     while (j < n2) arr[k++] = right[j++];
b50     return res;
35e }

462 int countInv(vector<int>& arr, int l, int r){
11e     int res = 0;
b14     if (l < r) {
ec9         int m = (r + l) / 2;
453         res += countInv(arr, l, m);
257         res += countInv(arr, m + 1, r);
470         res += countAndMerge(arr, l, m, r);
7e6     }
b50     return res;
7d1 }

92c int inversions(vector<int> &arr) {
851   	int n = arr.size();
cd9   	return countInv(arr, 0, n-1);
997 }
\end{lstlisting}

\subsection{Kadane}
\begin{lstlisting}
559 int kadane(const & vector<int> x) {
66c 	int sum = x[0];
fca 	int mx = x[0];
b19 	for (int i = 1; i < x.size(); i++) {
cbb 		sum = max(sum + x[i], x[i]);
2f2 		mx = max(mx, sum);
7b0 	}
b9a 	return mx;
df8 }

\end{lstlisting}

\subsection{Number of Pairs (GCD e LCM)}
\begin{lstlisting}
28b alguma coisa de manipular a expressao ja que o lcm(a,b) = a*b/gcd(a,b)
4e4 x = c*(a*b/gcd(a,b)) - d*gcd(a,b)
618 gcd(a,b) = g
 
fa2 lcm(a,b) = g*x*y onde gcd(x,y) = 1
 
12c por ex se x = 5 e y = 3
3d2 lcm(5,3) = 5*3*gcd(5,3){1} = 15
749 se a,b nao forem coprimos achamos um x,y que fazem eles serem coprimos, de maneira que
 
3c9 a = 4
82b b = 6
 
9f7 lcm(a,b) = 12
20a gcd(a,b) = 2
 
4f8 lcm(a,b) = x*y*gcd(a,b)
701 12 = x*y*2
12d x*y = 6
587 x = 3
983 y = 2
583 ou 
6d1 x = 6
6a2 y = 1
 
e36 lcm(4,6) = 3*2*2 correto
64f lcm(4,6) = 6*1*2 correto
 
578 entao lcm(a,b) = x*y*gcd(a,b)
 
672 x = c*x*y*g - d*g
b83 x = g((c*x*y) - d)
856 x/g = c*x*y - d
c4c */
 
 
f09 int spf[MAX], factors[MAX];
 
0a8 void build(){
81f     spf[1] = 1;
ce4     for (int i = 2; i < MAX; i+=2) spf[i] = 2;
7b9     for (int i = 3; i < MAX; i+=2){
e8b         if (spf[i] == 0){
cdc             spf[i] = i;
37e             for (int j = i; j*i < MAX; j+=2){
9fe                 if (spf[i*j] == 0) spf[i*j] = i;
2bf             }
859         }
8e9     }
c29 }
 
cbc void buildfactors(){
8db     for (int i = 2; i < MAX; i++){
4b7         factors[i] = factors[i/spf[i]];
b6b         if (spf[i/spf[i]] != spf[i]) factors[i]++;
2ba     } 
9cb }
 
c09 vector<int> getdiv(int x){
5ce     vector<int> d;
09a     for (int i = 1; i * i <= x; i++) {
72f         if (x % i == 0){ 
f9a             d.push_back(i);
331             if (i*i!=x) d.push_back(x/i);
7d0         }
a93     }
bae     sort(all(d));
be2     return d;
337 }
 
63d void solve(){
218     int c, d, x;
234     cin >> c >> d >> x;
0e8     vector<int> divisors = getdiv(x);
1a4     int ans = 0;
3fa     for (auto gc : divisors){
266         if (((x/gc)+d) % c != 0) continue;
ce4         int goal = (x/gc + d) / c;
            //goal significa que eu quero achar dois caras que multiplicados sejam goal, e o gcd deles seja 1
            //nisso eu vou descobrir a2 e b2, para achar a e b eu multiplico pelo u (gcd(a,b)) e calculo o lcm e depois testo os pares n sei como ainda
            //aqui eu acho o lcm e consigo descobrir o gcd de a e b
            // no caso o gcd ja eh meu divisor
            // ai basta achar as combinacoes validas de a e b
            // fatorar o goal e ver os fatores primos, para cada fator primo, a e b podem optar por receber ou nao ele
            // todos os fatores primos precisam ser usados, ou em a ou em b
            // entao temos como resposta a adicao de todas as combinacoes de k (qtd de fatores primos) em a,b
6fc         int k = factors[goal];
            //debug(k);
741         ans += (1<<k);   
495     }
886     cout << ans << endl;
450 }
 
0dd signed main(){
5d6     winton;
6f2     build();
d37     buildfactors();
8bd     int t;
f1b     cin >> t;
6e6     while(t--)solve();
af8 }
\end{lstlisting}

\subsection{Path Queries 2}
\begin{lstlisting}
2b7 #include <bits/stdc++.h>
ca4 using namespace std;
 
b63 const int MAX = 2e5+7;

fb1 namespace SegTree {
1a8     int n;
8f8     int pot2;
562     vector<int> tree;
    
0d8     void build(int n2, int* v2) {
1e3         n = n2;
fa5         pot2 = 1;
2c6         while (pot2 < n) pot2 <<= 1;
26a         tree.assign(2 * pot2, 0);
4e2         for (int i = 0; i < n; ++i) tree[pot2 + i] = v2[i];
813         for (int i = pot2 - 1; i >= 1; --i) tree[i] = max(tree[2*i], tree[2*i+1]);
641     }
    
4ea     int query(int a, int b) {
026         if (a > b) return 0;
238         a += pot2; b += pot2;
11e         int res = 0;
cce         while (a <= b) {
16c             if (a & 1) res = max(res, tree[a++]);
410             if (!(b & 1)) res = max(res, tree[b--]);
691             a >>= 1; b >>= 1;
80e         }
b50         return res;
d03     }
    
cf6     void update(int i, int x) {
1a7         int p = i + pot2;
3bc         tree[p] = x;
8b8         p >>= 1;
f77         while (p >= 1) {
280             tree[p] = max(tree[2*p], tree[2*p+1]);
8b8             p >>= 1;
6f4         }
4ad     }
092 }

a70 namespace HLD{
b24     vector<int> graph[MAX];
80f     int timer, h[MAX], ancestor[MAX], sz[MAX], pos[MAX], base[MAX], peso[MAX];
        
696     void build_hld(int v, int p = -1, int f = 1) {
c06         base[pos[v] = timer++] = peso[v]; 
588         sz[v] = 1;
504         for (auto &i : graph[v]) if (i != p) {
983             ancestor[i] = v;
71b             h[i] = (i == graph[v][0] ? h[v] : i);
e76             build_hld(i, v, f); sz[v] += sz[i];
    
c51             if (sz[i] > sz[graph[v][0]] or graph[v][0] == p) swap(i, graph[v][0]);
4e8         }
fd6         if (p*f == -1) build_hld(h[v] = v, -1, timer = 0);
bb3     }
    
1f8     void build(int root = 0) {
451         timer = 0;
295         build_hld(root);
160         SegTree::build(timer, base);
e95     }
    
f04     int query_path(int a, int b) {
11e         int res = 0;
9a1         while (h[a] != h[b]) {
aa1             if (pos[a] < pos[b]) swap(a, b);
30d             res = max(res, SegTree::query(pos[h[a]], pos[a]));
78b             a = ancestor[h[a]];
f86         }
9bc         if (pos[a] > pos[b]) swap(a, b);
bb7         res = max(res, SegTree::query(pos[a], pos[b]));
b50         return res;
852     }
        
081     void update_node(int i, int x) {
6e7         peso[i] = x;
2dc         SegTree::update(pos[i], x);
b94     }
95d }


0dd signed main(){
3b9     int n, q;
55a     cin >> n >> q;
04b     for (int i = 0; i < n; i++) cin >> HLD::peso[i];
6f5     for (int i = 1; i < n; i++){
ba2         int a, b;
0a8         cin >> a >> b;
e67         --a; --b;
770         HLD::graph[a].push_back(b);
ddb         HLD::graph[b].push_back(a);
364     }
b9a     HLD::build();
1f4     while(q--){
961         int type;
7b9         cin >> type;
821         if (type == 1){
8d2             int s, x;
f3c             cin >> s >> x;
351             s--;
847             HLD::update_node(s,x);
8e7         }
4e6         else {
ba2             int a, b;
0a8             cin >> a >> b;
450             a--; b--;
0ef             cout << HLD::query_path(a,b) << endl;
0ac         }
0a4     }
fba }
\end{lstlisting}

\subsection{Query de Segtree Offline}
\begin{lstlisting}
a03 int n, tree[4*MAX];

d42 int query (int node, int l, int r, int a, int b){
d5e     if (b < l || r < a) return 0;
14b     if (l >= a && r <= b)  return tree[node];
ee4     int m = (l+r)/2; 
5a3     return query(node*2, l, m, a, b) + query(node*2+1, m+1, r, a, b);
287 }

27f void update(int node, int l, int r, int i, int x){
791     if (i < l or i > r) return;
893     if (l == r) {
d6d         tree[node]=x;
505         return;
b74     }
ee4     int m = (l+r)/2;
3b2     update(node*2, l, m, i, x);
29a     update(node*2+1, m+1, r, i, x);
a16     tree[node] = tree[node*2] + tree[node*2+1];
d72 }

0dd signed main(){
5d6     winton;
45f     int q;
55a     cin >> n >> q;
f96     vector<pair<int,int>> a(n);
603     for (int i = 0; i < n; i++) {
8d6         cin >> a[i].first;
c92         a[i].second = i;
bd5     }
816     vector<tuple<int,int,int,int>> queries;
abf     for (int i = 0; i < q; i++){
12b         int l, r, m;
4c8         cin >> l >> m >> r;
0ef         if (a[l-1].first < m && r > 0){
205             queries.push_back ({m, l, l+r-1, i});
705         }
342     }
6a2     sort(rall(a));
7a3     sort(rall(queries));
bec     int pos = 0;
4cf     vector<int>ans(q,0);
5cc     for (auto [m,l,r,i] : queries){
730         while (pos < n && a[pos].first >= m){
                // debug(m);
                // debug(l);
                // debug(r);
                // debug(a[pos].second);
d85             update(1,0,n-1,a[pos].second,1);
657             pos++;
f08         }
a09         ans[i] = query(1,0,n-1,l,r);
            // cout << tree[8] << " " << tree[9] << " ";
            // for (int i = 5; i <= 7; i++)cout << tree[i] << " ";
            // cout << endl;
088     }
16b     for (auto u : ans) cout << u << endl;
e2b }
\end{lstlisting}

\subsection{Simspons Paradox}
\begin{lstlisting}
//Trend or pattern that appears in different groups of data is inconsistent 
//(disappears or even reverses) with what we see when the groups are combined.
//
//
//In this problem you are to construct an example of the dataset illustrating the Simpson's paradox.0
//More specifically, let us assume (similarly to the example above) that there are two teams who have been 
//solving problems of N categories and the total number of problems solved by each of the teams is M
//Let us denote the number of the problems in i-th category solved by Team X as ai, attempted by bi
//Similarly, let us define ci as the number of problems solved by Team Y in the i-th category and di as the number of problems attempted.

0dd signed main(){
14e     int n, m;
aa3     cin >> n >> m; 
        // a1 = b1 = 1, isso da um ratio de 1 usando so uma questao
        // c1 = d1-1, isso da um ratio de 0.99 usando quantas questoes eu quiser
        // so resta calcular para as restantes das categorias
        
8dd     int d1 = m - 3*(n-1);
        //debug(d1);
4d4     int c1 = d1-1;
93a     cout << 1 << " " << 1 << " " << c1 << " " << d1 << endl;
b28     for (int i = 0; i < n-2; i++){
0cd         cout << 1 << " " << 2 << " " << 1 << " " << 3 << endl;
159     }
abc     int bn = m-((n-2)*2)-1;
5da     int dn = m-d1-((n-2)*3);
816     int an = (c1 - 1);
c56     int cn = 1;
393     cout << an << " " << bn << " " << cn << " " << dn << endl;
cbd }
\end{lstlisting}

\subsection{Steiner Tree}
\begin{lstlisting}
// steiner: retorna o peso da menor arvore que cobre os vertices S
// get_steiner: retorna o valor minimo e as arestas de uma solucao
// se nao tiver solucao retorna LINF
//
// grafo nao pode ter pesos negativos
// se so tiver peso nas arestas/vertices pode deletar os vw/w no codigo
//
// k = |S|
// O(3^k * n + 2^k * m log m)

// otimizacao: joga um vertice x do S fora e pegue a resposta em dp[...][x] e reconstrua a arvore a partir dele
// ta comentado no codigo as mudancas necessarias

1a8 int n; // numero de vertices
c0d vector<pair<int, int>> g[MAX]; // {vizinho, peso}
920 ll d[1 << K][MAX]; // dp[mask][v] = arvore minima com o subconjunto mask de S e o vertice v
0b0 ll vw[MAX]; // peso do vertice

c8f ll steiner(const vector<int> &S) {
934 	int k = S.size(); // k--;
dea 	for (int mask = 0; mask < (1 << k); mask++) for(int v = 0; v < n; v++) d[mask][v] = LINF;
6b8 	for (int v = 0; v < n; v++) d[0][v] = vw[v];
254 	for (int i = 0; i < k; ++i) d[1 << i][S[i]] = vw[S[i]];
042 	for (int mask = 1; mask < (1 << k); mask++) {
b5b 		for (int a = (mask - 1) & mask; a; a = (a - 1) & mask) {
638 			int b = mask ^ a;
6bf 			if (b > a) break;
84d 			for (int v = 0; v < n; v++)
2e6 				d[mask][v] = min(d[mask][v], d[a][v] + d[b][v] - vw[v]);
2ab 		}
88c 		priority_queue<pair<ll, int>> pq;
778 		for (int v = 0; v < n; v++) {
6ad 			if (d[mask][v] == LINF) continue;
5ca 			pq.emplace(-d[mask][v], v);
f2e 		}
265 		while (pq.size()) {
a25 			auto [ndist, u] = pq.top(); pq.pop();
dad 			if (-ndist > d[mask][u]) continue;
c38 			for (auto [idx, w] : g[u]) if (d[mask][idx] > d[mask][u] + w + vw[idx]) {
679 				d[mask][idx] = d[mask][u] + w + vw[idx];
a2e 				pq.emplace(-d[mask][idx], idx);
07e 			}
de5 		}
a65 	}
478 	return d[(1 << k) - 1][S[0]]; // S[k]
        //para pegar o custo minimo para juntar todos os S's basta percorrer a mask-1 para todo no e ver qual o menor
704 }

d41 #warning se k=1 a solucao eh a folha isolada e a funcao retorna edg = {}
d41 #warning se k=0 crasha
4d1 pair<ll,vector<pair<int,int>>> get_steiner(const vector<int> &S) {
934 	int k = S.size(); // k--;
8ec 	ll ans = steiner(S);
c8d 	vector<pair<int,int>> edg;
57f 	stack<pair<int,int>> stk;
09b 	stk.emplace((1 << k) - 1, S[0]); // S[k]
07d 	while (!stk.empty()) {
c37 		bool cont = 0;
9c6 		auto [mask,u] = stk.top();stk.pop();
de2 		if ((__builtin_popcount(mask) == 1 and u == S[__bit_width(mask) - 1])) continue;
851 		for (auto [idx, w] : g[u]){
bb4 			if (d[mask][u] == d[mask][idx] + w + vw[u]) {
fc7 				edg.emplace_back(u, idx);
8ab 				stk.emplace(mask, idx);
a04 				cont = true;
c2b 				break;
342 			}
ed9 		}
3b5 		if (cont) continue;
b5b 		for (int a = (mask - 1) & mask; a; a = (a - 1) & mask) {
638 			int b = mask ^ a;
1e8 			if (d[mask][u] == d[a][u] + d[b][u] - vw[u]) {
be8 				stk.emplace(a, u);
c0a 				stk.emplace(b, u);
a04 				cont = true;
c2b 				break;
f52 			}
d29 		}
c5c 		assert(!mask || cont);
2b1 	}
be8 	return {ans, edg};
cf6 }
\end{lstlisting}

\subsection{Subtree Queries }
\begin{lstlisting}
5c1 const int MAXN = 2e5+7;

fff int tree[4*MAXN], tin[MAXN], tout[MAXN], flat[MAXN];
309 vector<vector<int>>graph;
2dc int timer = 0;

955 void dfs(int v, int p){
237     tin[v] = ++timer;
284     for (auto u : graph[v]){
58d         if (u!=p){
412             dfs(u, v);
bb2         }
ab6     }
a66     tout[v] = timer;
03d }

7a2 int build(int node, int l, int r){
893     if (l == r){
019         return tree[node] = flat[l];
f94     }
ee4     int m = (l+r)/2;
174     return tree[node] = build(node*2, l, m) + build(node*2+1, m+1, r);
97f }

d42 int query (int node, int l, int r, int a, int b){
d5e     if (b < l || r < a) return 0;
14b     if (l >= a && r <= b)  return tree[node];
ee4     int m = (l+r)/2; 
5a3     return query(node*2, l, m, a, b) + query(node*2+1, m+1, r, a, b);
287 }

27f void update(int node, int l, int r, int i, int x){
791     if (i < l or i > r) return;
893     if (l == r) {
d6d         tree[node]=x;
505         return;
b74     }
ee4     int m = (l+r)/2;
3b2     update(node*2, l, m, i, x);
29a     update(node*2+1, m+1, r, i, x);
a16     tree[node] = tree[node*2] + tree[node*2+1];
d72 }

0dd signed main(){
3b9     int n, q;
55a     cin >> n >> q;
979     vector<int>arr(n+1);
78a     for (int i = 1; i <= n; i++){
5d3         cin >> arr[i];
370     }
8f4     graph.resize(n+1);
6f5     for (int i = 1; i < n; i++){
ba2         int a, b;
0a8         cin >> a >> b;
8fa         graph[a].push_back(b);
4c6         graph[b].push_back(a);
080     }
b23     dfs(1,1);
78a     for (int i = 1; i <= n; i++){
a0b         flat[tin[i]] = arr[i];
616     }
be6     build(1,1,n);
        
1f4     while(q--){
ffd         int op;
f99         cin >> op;
e55         if (op == 1){
d5f             int idx, valor;
066             cin >> idx >> valor;
931             update(1, 1, n, tin[idx], valor);
fcf         }
4e6         else {
8ab             int s;
a18             cin >> s;
    
16b             int ans = query (1,1,n,tin[s], tout[s]);
886             cout << ans << endl;
97e         }
            
d0f     }
a7b }
\end{lstlisting}



%%%%%%%%%%%%%%%%%%%%
%
% Matemática
%
%%%%%%%%%%%%%%%%%%%%

\section{Matemática}

\subsection{Combinatoria}
\begin{lstlisting}
3e4 const int MOD = 1e9+7;
4e9 const int MOD = 998244353;

183 const int MAX = 1e6+7;

e4a int fat[MAX], invfat[MAX];

bf9 int fastExpo(int base, int exp) {
c92     int res = 1;
fb9     while(exp) {
8ca         if (exp & 1) res = res * base % MOD;
a6a         base = base * base % MOD;
ef0         exp >>= 1;
635     }
542     return res%MOD;
2e3 }

868 int modinv(int n) {
562     return fastExpo(n, MOD - 2);
92e }

e9d int modiv(int a, int b) {
a7c     return (a * modinv(b)) % MOD;
a02 }

0a8 void build(){
557     fat[0] = 1;
38a     for (int i = 1; i < MAX; i++) fat[i] = (fat[i-1] * i) % MOD;
507     invfat[MAX-1] = modinv(fat[MAX-1]);
e92     for (int i = MAX - 2; i >= 0; i--) invfat[i] = (invfat[i+1] * (i+1)) % MOD;
4fa }
c52 int comb(int a, int b){
a75     if (b > a || b < 0 || a < 0) return 0;
97d     return (((fat[a] * invfat[b]) % MOD) * invfat[a-b]) % MOD;
395 }

76f int combinations(int n, int k) {
        // Basic edge cases
032     if (k < 0 || k > n) {
bb3         return 0;
edf     }
b62     if (k == 0 || k == n) {
6a5         return 1;
f14     }
dd2     if (k > n - k) {
1b5         k = n - k;
b13     }
c92     int res = 1;
64b     for (int i = 1; i <= k; ++i) {
e56         res = res * (n - i + 1) / i;
ad4     }
b50     return res;
ff6 }

34a int combinations_mod(int n, int k) {
e29     if (k < 0 || k > n)  return 0;
bbe     if (k == 0 || k == n) return 1;
63f     if (k > n / 2) k = n - k;
bd9     int numerator = 1;
599     for(int i = 0; i < k; i++) numerator = (numerator * (n - i)) % MOD;
cd4     int denominator = 1;
e49     for(int i = 1; i <= k; i++) denominator = (denominator * i) % MOD;
931     return (numerator * modinv(denominator)) % MOD;
623 }
\end{lstlisting}

\subsection{ctor<int> divisors;}
\begin{lstlisting}
b32 if (tot % i == 0) {
025     divisors.push_back(i);
d73     if (i * i != tot) divisors.push_back(tot / i);
7b3 }
\end{lstlisting}

\subsection{EV Máximo}
\begin{lstlisting}
/*
799 (r/m)^n - ((r-1)/m)^n = probabilidade de ter exatamente r como o maior resultado 
b85 somatorio de P(r) * r = EV do maximo

e08 n: numero de eventos
9cb r: resultado
0ee m: quantidade de resultados possiveis para um evento
c4c */

0e8 r * ((fastExpo((r/m), n)) - (fastExpo((ld)(r-1)/m, n)));
\end{lstlisting}

\subsection{Exponenciação Rápida}
\begin{lstlisting}
bf9 int fastExpo(int base, int exp) {
c92     int res = 1;
fb9     while(exp) {
8ca         if (exp & 1) res = res * base % MOD;
a6a         base = base * base % MOD;
ef0         exp >>= 1;
635     }
542     return res%MOD;
2e3 }
e9d int modiv(int a, int b){
659     return (((a % MOD )* (fastExpo(b, MOD-2) % MOD)) % MOD);
e50 }

bf9 int fastExpo(int base, int exp) {
c92     int res = 1;
fb9     while(exp) {
13e         if (exp & 1) res = res * base;
6c1         base = base * base;
ef0         exp >>= 1;
319     }
b50     return res;
4fe }

83b int fastExpo(int base, int exp, int m) {
c92     int res = 1;
054     base %= m;
fb9     while(exp) {
970         if (exp & 1) res = res * base % m;
103         base = base * base % m;
ef0         exp >>= 1;
268     }
17a     return res%m;
1b8 }

ff4 ld fastExpo(ld base, int exp) {
9d0     ld res = 1;
fb9     while(exp) {
13e         if (exp & 1) res = res * base;
6c1         base = base * base;
ef0         exp >>= 1;
319     }
b50     return res;
b1b }

\end{lstlisting}

\subsection{Fatoração Prima}
\begin{lstlisting}
4ae vector<int> factor(int n) {
fd3     vector<int> primes;
d91     for (int p = 2; p * p <= n; ++p) {
03e         while (n % p == 0) {
bce             primes.push_back(p);
f4a             n /= p;
6e7         }
4ac     }
f3d     if (n > 1) 
ccc         primes.push_back(n);
a20     return primes;
8de }


\end{lstlisting}

\subsection{GCD e LCM}
\begin{lstlisting}
ba6 int gcd(int a, int b) {
d65     return b == 0 ? a : gcd(b, a % b);
4c5 }

3ec int mmc(int a, int b) {
8b8     return a / gcd(a, b) * b;
4c5 }
\end{lstlisting}

\subsection{Phi Function}
\begin{lstlisting}
//quantos numeros sao coprimos com N de 1 ate N
633 int phi[MAX];

0ae void precomputePHI(){
104     for (int i = 0; i < MAX; i++)
ccf         phi[i] = i;
    
8db     for (int i = 2; i < MAX; i++) {
e32         if (phi[i] == i) {
d87             for (int j = i; j < MAX; j += i)
a9b                 phi[j] -= phi[j] / i;
0cc         }
f9d     }
159 }
\end{lstlisting}

\subsection{Segmented Sieve}
\begin{lstlisting}
576 vector<char> segmentedSieve(long long L, long long R) {
        // generate all primes up to sqrt(R)
4b2     long long lim = sqrt(R);
537     vector<char> mark(lim + 1, false);
d2d     vector<long long> primes;
    
b74     for (long long i = 2; i <= lim; ++i) {
c9d         if (!mark[i]) {
d38             primes.emplace_back(i);
1e8             for (long long j = i * i; j <= lim; j += i)
148                 mark[j] = true;
181         }
920     }
    
f9d     vector<char> isPrime(R - L + 1, true);
867     for (long long i : primes){
f8c         for (long long j = max(i * i, (L + i - 1) / i * i); j <= R; j += i){       
0d5             isPrime[j - L] = false;
cd3         }
24a     }
86a     if (L == 1) isPrime[0] = false;
        
        //para pegar os primos basta fazer i + L no vetor de isprime 
3c8     return isPrime;
a43 }
\end{lstlisting}

\subsection{Sieve}
\begin{lstlisting}
fd3 vector<int> primes;

8e0 void sieve() {
a94     vector<bool> is_prime(MAX + 1, true);
428     is_prime[0] = false;
8fc     is_prime[1] = false;
31f     for (int p = 2; p * p <= MAX; ++p) {
29c         if (is_prime[p]) {
305             for (int multiple = p * p; multiple <= MAX; multiple += p) {
4f8                 is_prime[multiple] = false;
16d             }
354         }
d56     }
e65     for (int i = 2; i <= MAX; ++i) {
4a3         if (is_prime[i]) {
e74             primes.push_back(i);
34c         }
b9d     }
ee0 }
\end{lstlisting}

\subsection{Smallest Prime Factor}
\begin{lstlisting}
b32 int n, spf[MAX];

0a8 void build(){
ce4     for (int i = 2; i < MAX; i+=2) spf[i] = 2;
7b9     for (int i = 3; i < MAX; i+=2){
e8b         if (spf[i] == 0){
cdc             spf[i] = i;
37e             for (int j = i; j*i < MAX; j+=2){
9fe                 if (spf[i*j] == 0) spf[i*j] = i;
2bf             }
859         }
8e9     }
5ce }

826 vector<int> factorize(int x){
fd3     vector<int> primes;
40e     while (x > 1){
80a         int p = spf[x];
bce         primes.push_back(p);
15a         while(x%p == 0) x/=p;
64d     }
        
a20     return primes;
0f8 }

//Numero de fatores primos diferentes que fatoram cada numero
11f int spf[MAX], pf[MAX];

0a8 void build(){
81f     spf[1] = 1;
ce4     for (int i = 2; i < MAX; i+=2) spf[i] = 2;
7b9     for (int i = 3; i < MAX; i+=2){
e8b         if (spf[i] == 0){
cdc             spf[i] = i;
37e             for (int j = i; j*i < MAX; j+=2){
9fe                 if (spf[i*j] == 0) spf[i*j] = i;
2bf             }
859         }
8e9     }
c29 }

da9 void buildspf(){
8db     for (int i = 2; i < MAX; i++){
e09         pf[i] = pf[i/spf[i]];
a28         if (spf[i/spf[i]] != spf[i]) pf[i]++;
3a0     } 
d85 }
\end{lstlisting}



%%%%%%%%%%%%%%%%%%%%
%
% DP
%
%%%%%%%%%%%%%%%%%%%%

\section{DP}

\subsection{==========================================}
\begin{lstlisting}
// ==========================================
730 int dp[20][11][2][2]; // [index][state][tight][ldz]

c3f int pd(int index, int state, int tight, int ldz) {
2e8     if(index == n) {
505         return /* condicao base */;
d0c     }
        
890     if (dp[index][state][tight][ldz] != -1) return dp[index][state][tight][ldz];
    
c23     int ub = tight ? r[index] - '0' : 9;
1a4     int ans = 0;
    
70b     for (int digit = 0; digit <= ub; digit++) {
            // if(!ldz && condicao_invalida) continue;
    
e16         int new_tight = tight && (digit == ub);
d3b         int new_ldz = ldz && (digit == 0);
            
f1b         int new_state = state; // atualiza seu estado
    
cf0         ans += pd(index + 1, new_state, new_tight, new_ldz);
572     }
c77     return dp[index][state][tight][ldz] = ans; 
18b }

// ==========================================
// TEMPLATE 2: RECURSIVA Range [L, R] Simultâneo
// ==========================================
377 int dp[20][11][2][2][2]; // [index][state][above][under][ldz]

d6c int pd(int index, int state, int above, int under, int ldz) {
817     if(index == n) return /* condicao base */;
    
2d4     if (dp[index][state][above][under][ldz] != -1) return dp[index][state][above][under][ldz];
    
c62     int ub = above ? r[index] - '0' : 9;
f28     int lb = under ? l[index] - '0' : 0;
1a4     int ans = 0;
    
369     for (int digit = lb; digit <= ub; digit++) {
            // if(!ldz && condicao_invalida) continue;
    
394         int new_above = above && (digit == ub);
4ff         int new_under = under && (digit == lb);
d3b         int new_ldz = ldz && (digit == 0);
            
f1b         int new_state = state; // atualiza seu estado
    
a39         ans += pd(index + 1, new_state, new_above, new_under, new_ldz);
510     }
430     return dp[index][state][above][under][ldz] = ans; 
f74 }

// ==========================================
// TEMPLATE 3: ITERATIVA PUSH DP
// ==========================================
518 void solve_iterative() {
730     int dp[20][11][2][2]; // [index][last][tight][ldz]
4b9     memset(dp, 0, sizeof(dp));
    
a87     dp[0][10][1][1] = 1;
    
03b     for (int index = 0; index < n; index++) {
0c8         for (int last = 0; last <= 10; last++) {
cc3             for (int tight = 0; tight < 2; tight++) {
35d                 for (int ldz = 0; ldz < 2; ldz++) {
                        
dcf                     int ways = dp[index][last][tight][ldz];
e6d                     if(!ways) continue;
    
b57                     int limit = tight ? r[index] - '0' : 9;
                        
865                     for (int digit = 0; digit <= limit; digit++) {
                            
5c6                         if(!ldz && last == digit) continue;
    
14f                         int new_tight = tight && (digit == limit);
d3b                         int new_ldz = ldz && (digit == 0);
                            
ae9                         int next = new_ldz ? 10 : digit;
    
b2e                         dp[index+1][next][new_tight][new_ldz] += ways;
660                     }
cde                 }
32b             }
f27         }
35b     }
77e }


7d8 int pd(int index, int sum, int above, int under){
1df     if(index == n) return (sum == 0);
    
a73     if (!above && !under && dp[index][sum] != -1) return dp[index][sum];
    
c62     int ub = above ? r[index] - '0' : 9;
f28     int lb = under ? l[index] - '0' : 0;
1a4     int ans = 0;
    
369     for (int digit = lb; digit <= ub; digit++){
    
394         int new_above = above && (digit == ub);
4ff         int new_under = under && (digit == lb);
            
f98         int new_sum = (sum * 10 + digit) % M;
            
e79         ans = (ans + pd(index + 1, new_sum, new_above, new_under)) % MOD;
2ff     }
        
6aa     if (!above && !under) dp[index][sum] = ans; 
ba7     return ans;
bec }


//so precisa fazer memset uma vez mas dai precisa fazer padding pra ambos os numeros terem o mesmo numero de digitos
5db int count(int x){
d06     s = to_string(x);
        
eda     int pad = 19 - s.size();
a0e     if (pad > 0) s = string(pad, '0') + s;
        
25a     n = s.size();
272     return pd(0,0,1);
2b8 }

63d void solve(){
74f     int l, r;
5a9     cin >> l >> r;
27a     memset(dp, -1, sizeof(dp));
aba     cout << count(r) - count(l-1) << endl;
21c }

/*
b73 otimizacao super importante
7ea desse jeito, o digito menos significativo sempre esta na posicao 0
bb3 entao as posicoes dos digitos de cada numero sempre estarao na mesma posicao
2c3 por exemplo 1234 e 12
c6f do jeito normal ficaria 
81d 1234
c20   12

47a desse jeito fica
d93 4321
3c5   21

4be entao a potencia de 10 em cada posicao e sempre a mesma, permitindo so rodar o memset uma vez 
c4c */
990 vector<int> v; 
7c2 int pd(int index, int tight, int mask, int ldz){
789     if (index == -1) return ((__bit_width(mask) - 1) == __popcount(mask));
        
090     if (!tight && dp[index][mask][ldz] != -1) return dp[index][mask][ldz];
    
cce     int ub = tight ? v[index] : 9;
1a4     int ans = 0;
a93     int sum = 0;
    
70b     for (int digit = 0; digit <= ub; digit++){
e16         int new_tight = tight && (digit == ub);
    
d3b         int new_ldz = ldz && (digit == 0);
            
963         int new_mask = mask;
889         if (!new_ldz) new_mask |= (1 << digit);
    
6d6         ans += pd(index - 1, new_tight, new_mask, new_ldz);
739     }
    
0ab     if (!tight) dp[index][mask][ldz] = ans; 
ba7     return ans;
1c6 }

a7e int count(ll x){
848     v.clear();
f4c     while(x) {
d6f         v.push_back(x % 10);
6bd         x /= 10;
053     }
    
            /* [L,R]
3cd     int count(int x, int y){
d65     v2.clear();
f4c     while(x) {
b64         v2.push_back(x % 10);
6bd         x /= 10;
d12     }
848     v.clear();
    
1b8     while(y) {
da8         v.push_back(y % 10);
a53         y /= 10;
38c     }
de6     while(v2.size() < v.size()) v2.push_back(0LL);
c4c     */
       
d6c     n = v.size();
726     return pd(n-1,1,0,1);
6ae }
    
63d void solve(){
bb3     ll r;
da1     cin >> r;
226     cout << count(r) << endl;
c01 }
    
0dd signed main(){
5d6     winton;
27a     memset(dp, -1, sizeof(dp));
8bd     int t;
f1b     cin >> t;
6e6     while(t--) solve();
f37 }
\end{lstlisting}

\subsection{Edit Distance}
\begin{lstlisting}

ac1 string s, t;
fee cin >> s >> t;
3e6 int n = s.length();
32a int m = t.length();
78a for (int i = 1; i <= n; i++){
cbc     for (int j = 1; j <= m; j++){
c70         if(s[i-1] == t[j-1]) dp[i][j] = 1 + dp[i-1][j-1];
6d7         else dp[i][j] = max(dp[i-1][j], dp[i][j-1]);
9e1     }
4b0 }
\end{lstlisting}

\subsection{LCS}
\begin{lstlisting}

ac1 string s, t;
fee cin >> s >> t;
3e6 int n = s.length();
32a int m = t.length();
78a for (int i = 1; i <= n; i++){
cbc     for (int j = 1; j <= m; j++){
c70         if(s[i-1] == t[j-1]) dp[i][j] = 1 + dp[i-1][j-1];
6d7         else dp[i][j] = max(dp[i-1][j], dp[i][j-1]);
9e1     }
4b0 }
\end{lstlisting}



%%%%%%%%%%%%%%%%%%%%
%
% Grafo
%
%%%%%%%%%%%%%%%%%%%%

\section{Grafo}

\subsection{Bridge Tree}
\begin{lstlisting}

9fa int t, c, low[MAX], pre[MAX], visited[MAX], comp[MAX];
7b4 stack<int> st;
375 vector<int> graph[MAX], compgraph[MAX];

955 void dfs(int v, int p){
cda     visited[v] = 1;
98b     low[v] = pre[v] = ++t;
30c     st.emplace(v);
284     for (auto u : graph[v]){
264         if (!visited[u]){
412             dfs(u,v);
c80             low[v] = min(low[v], low[u]);
77d         }   
4e6         else {
2da             if (u == p) continue;
39b             low[v] = min(low[v], pre[u]);
acd         }
11d     }
cf1     if (low[v] == pre[v]){ //achamos uma bridge
6a4         c++;
2f4         int cur;
016         do {
596             cur = st.top();
25a             st.pop();
524             comp[cur] = c;
cca             compgraph[c].push_back(cur);
583         } while(cur != v);
b04     }
6fc }
\end{lstlisting}

\subsection{Djikstra}
\begin{lstlisting}
861 int dist[MAX]
d9c void djikstra(int st){
fe1     priority_queue<pair<int,int>> pq;
260     for (int i = 0; i < n; i++)dist[i]=LLONG_MAX;
d89     dist[st] = 0;
56c     pq.push({0,st});
502     while(!pq.empty()){
115         auto [d, v] = pq.top();
842         d = -d;
716         pq.pop();
dd8         if (d != dist[v]) continue;
940         for (auto [u,w] : graph[v]){
ac8             int  novadistancia = d + w;
bd8             if (novadistancia < dist[u]) {
d96                 dist[u] = novadistancia;
030                 pq.push({-novadistancia, u});
c98             }
379         }
133     }
6f3 }
\end{lstlisting}

\subsection{DSU }
\begin{lstlisting}
d56 struct DSU{
1a8     int n;
7d2     vector<int> p, r;
5c7     void init(int _n){
8fd         n = _n;
1c4         p.resize(n);
7e4         r.resize(n);
603         for(int i = 0; i < n; i++){
f13             p[i] = i;
fd1             r[i] = 0;
606         }
55e     }
    
57b     int find(int x){
c8c         return p[x] == x ? x : p[x] = find(p[x]);
2d3     }
    
765     bool unite(int a, int b){
bca         a = find(a);
b88         b = find(b);
5c2         if(a == b) return false;
dac         if(r[a] == r[b]) r[a]++;
d1b         if(r[a] > r[b]){
264             p[b] = a;
4e6             r[a]++;
db5         }else{
84b             p[a] = b;
f6c             r[b]++;
c3d         }
8a6         return true;
632     }
31c }


d56 struct DSU {
825 	vector<int> id, sz;
    
a56 	DSU(int n) : id(n), sz(n, 1) { iota(id.begin(), id.end(), 0); }
    
0cf 	int find(int a) { return a == id[a] ? a : id[a] = find(id[a]); }
    
440 	void unite(int a, int b) {
605 		a = find(a), b = find(b);
d54 		if (a == b) return;
956 		if (sz[a] < sz[b]) swap(a, b);
6d0 		sz[a] += sz[b], id[b] = a;
ea7 	}
c1b };

\end{lstlisting}

\subsection{HLD - aresta}
\begin{lstlisting}
// SegTree de soma
// query / update de soma das arestas
//
// Complexidades:
// build - O(n)
// query_path - O(log^2 (n))
// update_path - O(log^2 (n))
// query_subtree - O(log(n))
// update_subtree - O(log(n))

// namespace SegTree { ... }

a70 namespace HLD {
af8     vector<pair<int,int>> graph[MAX];
480     int timer, pos[MAX], sz[MAX], sobe[MAX], ancestor[MAX], h[MAX], v[MAX]; 
    
c2a     void build_hld(int node, int p = -1, int f = 1){
f20         v[pos[node] = timer++] = sobe[node];
57c         sz[node] = 1;
50f         for (auto &i : graph[node]){
dd2             auto [u, w] = i;
2da             if (u == p) continue;
828             sobe[u] = w;
df7             ancestor[u] = node;
112             h[u] = (i == graph[node][0] ? h[node] : u); 
b79             build_hld(u, node, f);
e0e             sz[node] += sz[u];
    
996             if (sz[u] > sz[graph[node][0].first] || graph[node][0].first == p){
8c8                 swap(i, graph[node][0]);
ee5             }
2b5         }
469         if (p*f == -1)build_hld(h[node] = node, -1, timer = 0);
0a5     }
    
1f8     void build(int root = 0){
451         timer = 0;
295         build_hld(root);
cbf         SegTree::build(timer, v);
517     }
    
f04     int query_path(int a, int b){
2d5         if (a == b) return 0;
aa1         if (pos[a] < pos[b]) swap(a,b);
186         if (h[a] == h[b]) return SegTree::query(pos[b] + 1, pos[a]);
587         return SegTree::query(pos[h[a]], pos[a]) + query_path(ancestor[h[a]], b);
cd0     }
    
920     void update_path(int a, int b, int x){
d54         if (a == b) return;
aa1         if (pos[a] < pos[b]) swap(a,b);
    
1eb         if (h[a] == h[b]) return (void)SegTree::update(pos[b]+1, pos[a], x);
046         SegTree::update(pos[h[a]], pos[a], x);
ab0         update_path(ancestor[h[a]], b, x);
10e     }
    
3e4     int query_subtree(int a){
b9f         if(sz[a] == 1) return 0;
f97         return SegTree::query(pos[a]+1, pos[a] + sz[a]-1);
ce2     }
    
acc     void update_subtree(int a, int x){
a5a         if(sz[a] == 1) return;
84f         SegTree::update(pos[a]+1, pos[a] + sz[a]-1, x);
bf2     }
    
7be     int lca(int a, int b){
aa1         if (pos[a] < pos[b]) swap(a,b);
0f9         return h[a] == h[b] ? b : lca(ancestor[h[a]], b);
263     }
b5d     int dist(int a, int b) {
2d5         if (a == b) return 0;
aa1         if (pos[a] < pos[b]) swap(a,b);
30a         if (h[a] == h[b]) return pos[a] - pos[b];
13e         return pos[a] - pos[h[a]] + 1 + dist(ancestor[h[a]], b);
58d     }
16d }
\end{lstlisting}

\subsection{Kosaraju}
\begin{lstlisting}
881 int n, com;
042 vector<int> g[MAX];
58d vector<int> gi[MAX];
981 vector<int> g_scc[MAX]; 
c5a int vis[MAX];
ee6 stack<int> S;
a52 int comp[MAX];

1ca void dfs(int k) {
59a 	vis[k] = 1;
7b1 	for (int i = 0; i < g[k].size(); i++)
8d5 		if (!vis[g[k][i]]) dfs(g[k][i]);
    
58f 	S.push(k);
0a1 }

436 void scc(int k, int c) {
59a 	vis[k] = 1;
52c 	comp[k] = c;
309 	for (int i = 0; i < gi[k].size(); i++)
bf6 		if (!vis[gi[k][i]]) scc(gi[k][i], c);
b43 }

db8 void kosaraju() {
991 	for (int i = 0; i < n; i++) vis[i] = 0;
    
158 	for (int i = 0; i < n; i++) if (!vis[i]) dfs(i);
    
991 	for (int i = 0; i < n; i++) vis[i] = 0;
d32 	while (S.size()) {
70b 		int u = S.top();
7de 		S.pop();
1d6 		if (!vis[u]) {
fd4             com++;
7d8             scc(u, com);
35b         }
96d     }
f84 }

// funcao para construir o grafo de componentes
ad1 void build_component_graph() {
687     for (int u = 0; u < n; u++) {
4d5         for (int v : g[u]) {
1ec             if (comp[u] != comp[v]) {
6d1                 g_scc[comp[u]].push_back(comp[v]);
5d9             }
b1e         }
1fb     }
        //remove duplicatas
002     for (int i = 1; i <= com; i++){
0a2         sort(g_scc[i].begin(), g_scc[i].end());
1f9         if (g_scc[i].empty())continue;
4ad         g_scc[i].erase(unique(g_scc[i].begin(), g_scc[i].end()), g_scc[i].end());
1d4     }
6b7 }
\end{lstlisting}

\subsection{LCA Binary Lifting}
\begin{lstlisting}
039 int n, ancestor[MAXN][LOG+1], depth[MAXN], visited[MAXN];
309 vector<vector<int>> graph;

955 void dfs(int v, int p){
cda     visited[v] = 1;
d38     ancestor[v][0] = p;
    
c64     for (int i = 1; i <= LOG; i++){
d14         ancestor[v][i] = ancestor[ancestor[v][i-1]][i-1];
a6e     }
    
284     for (auto u : graph[v]){
264         if (!visited[u]){
cf4             depth[u] = depth[v]+1;
412             dfs(u, v);
813         }
2a4     }
f1b }

c5d int lca(int v, int u){
6ed     if (depth[v] < depth[u]) swap(v,u);
    
4f1     for (int i = LOG-1; i >=  0; i--){
01e         if ((depth[v] - (1<<i)) >= depth[u]){
ed6             v = ancestor[v][i];
cc6         } 
830     }
    
d77     if (u == v) return v;
    
4f1     for (int i = LOG-1; i >= 0; i--){
c47         if (ancestor[v][i] != ancestor[u][i]){
ed6             v = ancestor[v][i];
fb0             u = ancestor[u][i];
329         }
921     }
6c2     return ancestor[v][0];
40e }

33e struct LCA {
        //one indexed
1c4     int n, l, cur;
47c     vector<vector<int>> adj, up;
991     vector<int> h, entry;
8c5     LCA (int n_) {
e95         n = n_;
4e0         l = 31 - __builtin_clz(n);
839         adj.resize(n + 1);
a71         up.assign(n + 1, vector<int>(l + 1));
717         h.resize(n + 1); entry.resize(n + 1);
b59         cur = 0;
521     }
50f     void add(int a, int b) {
fab         adj[a].push_back(b);
c87         adj[b].push_back(a);
386     }
17c     void dfs(int v = 1, int p = 0, int hh = 0) {
176         entry[v] = cur++;
4de         up[v][0] = p;
910         for(int i = 1; i <= l; i++) {
73a             up[v][i] = up[up[v][i - 1]][i - 1];
d67         }
9bc         for(auto u : adj[v]) {
58d             if(u != p) {
828                 dfs(u, v, hh + 1);
d5f             }
693         }
3bc         h[v] = hh;
a4d     }
7be     int lca(int a, int b) {
a6b         if(h[a] > h[b]) {
257             swap(a, b);
79b         }
14e         int dff = h[b] - h[a], cnt = 0;
d84         while(dff) {
57f             if(dff & 1) {
5be                 b = up[b][cnt];
0ce             }
f65             cnt++;
5e9             dff >>= 1;
488         }
a89         if(a == b) {
3f5             return a;
06e         }
957         for(int i = l; i >= 0; i--) {
c84             if(up[a][i] != up[b][i]) {
3f4                 a = up[a][i];
6c0                 b = up[b][i];
568             }
dc1         }
e6f         return up[a][0];
cb7     }
699     int raise(int node, int dist) {
            //if it returns 0, then it fucked up
ac9         int cnt = 0;
efd         while(dist) {
e3d             if(dist & 1) {
466                 node = up[node][cnt];
0ed             }
f65             cnt++;
705             dist >>= 1;
1c0         }
192         return node;
887     }
b5d     int dist(int a, int b) {
b80         int z = lca(a, b);
421         return h[a] + h[b] - 2 * h[z];
0a0     }
434 };
\end{lstlisting}

\subsection{LCA Sparse Table}
\begin{lstlisting}
/*
172 Fazer um euler tour na arvore e inserir todas as vezes que visitamos um no
6e7 Junto disso inserir a depth desse no em outro vetor
9f5 no final, o intervalo entre a e b sera representado por todos os nos que visitamos no caminho entre a e b
55b basta pegar o o no com altura minima nesse intervalo
31f para saber o ponto incial desse intervalo, basta pegar o tin de a e b

6e5 build O(N LOG N)
8fc query O(1)

b0d posso pensar em otimizar como chamo o lca guardando a ultima vez que o no foi chamado
c4c */

4f6 int n, timer = 0, tin[MAX], depth[MAX];
ba0 vector<int> et;
227 pair<int,int> sp[MAX*2][LOG+1];
309 vector<vector<int>> graph; //0 indexado

955 void dfs(int v, int p){
8c5     tin[v] = et.size();
490     et.push_back(v);
284     for (auto u : graph[v]){
58d         if (u!=p){
cf4             depth[u] = depth[v]+1;
412             dfs(u,v);
490             et.push_back(v);
f82         }
7ac     }
6e0 }

07b void buildtable(){
96d     for (int i = 0; i < (int)et.size(); i++){
1e7         sp[i][0] = {depth[et[i]], et[i]};
bcc     }
2ae     for (int j = 1; j <= LOG; j++){
538         for (int i = 0; i + (1<<(j-1)) < (int)et.size(); i++){
d30             sp[i][j] = min (sp[i][j-1], sp[i+(1<<(j-1))][j-1]);
985         }
fd8     }
641 }

2b4 pair<int,int> query (int a, int b){
bfb     int len =  b - a + 1;
5e4     int lg = 31 - __builtin_clz(len);
a57     return min(sp[a][lg], sp[b - (1<<lg) + 1][lg]);
11c }

7be int lca(int a, int b){
fb2     if (tin[a] > tin[b]) swap(a,b);
ce1     return query(tin[a], tin[b]).second;
940 }

0dd signed main() {
5d6     winton;
45f     int q;
55a     cin >> n >> q;
c02     graph.resize(n);
f50     for (int i = 0; i < n-1; i++){
ba2         int a, b;
0a8         cin >> a >> b;
9c7         graph[--a].push_back(--b);
4c6         graph[b].push_back(a);
601     }
564     dfs(0,0);
df6     buildtable();
1f4     while(q--){
ba2         int a, b;
0a8         cin >> a >> b;
450         a--;b--;
092         int mac = lca(a,b);
389         cout << depth[a] + depth[b] - 2*depth[mac] << endl;
e88     }
    
1e3 }
\end{lstlisting}

\subsection{Minimum Spanning Tree}
\begin{lstlisting}
d56 struct DSU {
8f3 	vector<int> comp, sz;
    
c1b 	DSU(int n) : comp(n), sz(n, 1) {
f7a 		iota(all(comp), 0);
576 	}
    
151 	int find(int u) {
dc3 		if (u == comp[u]) return u;
a07 		return comp[u] = find(comp[u]);
65f 	}
    
8bf 	bool merge(int a, int b) {
605 		a = find(a), b = find(b);
5c2 		if (a == b) return false;
afa 		sz[a] += sz[b];
ca8 		comp[b] = a;
8a6 		return true;
6bc 	}
33f };
 
63d void solve(){
14e     int n, m;
aa3     cin >> n >> m;
73b     vector<tuple<int,int,int>> edges; 
370     DSU dsu(n);
8f8     edges.resize(m);
dd5        for (int i = 0; i < m; i++){
684         int a, b, w;
060         cin >> a >> b >> w;
0d4         edges[i] = {w,a,b};
261     }
5c1     sort(all(edges));
704     int cost = 0;
dd5     for(int i = 0; i < m; i++){
13d         auto [w,u,v] = edges[i];
d19         if (dsu.merge(u,v)){
45f             cost += w;
b64         }
ce9     }
16f }
\end{lstlisting}

\subsection{Pontos de Articulação}
\begin{lstlisting}
/*
a7e Acha os articulation points do grafo
c4c */

f7d int t, c, low[MAX], pre[MAX], visited[MAX];
864 vector<int> graph[MAX], ans;

955 void dfs(int v, int p){
cda     visited[v] = 1;
98b     low[v] = pre[v] = ++t;
14e     bool art = false;
236     int filhos = 0;
284     for (auto u : graph[v]){
264         if (!visited[u]){
701             filhos++;
412             dfs(u,v);
c80             low[v] = min(low[v], low[u]);
    
b71             if(low[v] >= pre[v]) art = true;
4d8         }   
4e6         else {
2da             if (u == p) continue;
39b             low[v] = min(low[v], pre[u]);
acd         }
f4f     }
ca1     if (v == 1 && filhos >= 2) ans.push_back(v);
021     if (v != 1 && art) ans.push_back(v);
7fd }
\end{lstlisting}

\subsection{Toposort}
\begin{lstlisting}
// Fazer dfs pra pegar o indeg se necessario
0ad int indeg[MAX];

55b vector<int> order;
26a queue<int> q;

519 for (int i = 0; i < MAX; i++){
ac5     if (indeg[i] == 0)q.push(i);
3fd }
14d while(!q.empty()){
69f     int cur = q.front();   
ba5     order.push_back(cur);
833     q.pop();
411     for (auto u : graph[cur]){
472         indeg[u]--;
f90         if (indeg[u] == 0)q.push(u);
014     }
f6a }
\end{lstlisting}

\pagebreak


%%%%%%%%%%%%%%%%%%%%
%
% Extra
%
%%%%%%%%%%%%%%%%%%%%

\section{Extra}

\subsection{compressao.cpp}
\begin{lstlisting}
//compressao de coordenada para vector

vector<int> v(n);
vector<int> todos;
for (int i = 0; i < n; i++){
    cin >> v[i];
    todos.push_back(v[i]);
}

sort(all(todos));
todos.erase(unique(all(todos)), todos.end());
int m = todos.size();

auto comp = [&](int x) {
    return (lower_bound(all(todos), x) - todos.begin());
};

vector<int> comp(n);
for (int i = 0; i < n; i++){
    comp[i] = lower_bound(all(todos), v[i]) - todos.begin();
}
\end{lstlisting}

\subsection{lambda.cpp}
\begin{lstlisting}
long long ans = 0;

auto lambda_padrao = [&](int x, int y) -> int {
    return x + y;
};

auto lambda_generica = [&](const auto &v) {
    return (int)v.size();
};

auto lambda_recursiva = [&](auto &&self, int x) -> int {
    if (x <= 1) return 1;
    return x * self(self, x - 1);
};

function<void(int, int)> dfs = [&](int v, int p) {
    for (auto u :graph[v]){
        if (u == v) continue;
        dfs(u,v);
    }
};

auto reduce = [&](const multiset<string> &ms) -> multiset<int> {
    multiset<int> mi;
    for (const auto &s : ms) {
        if (s.size() > 1) {
            mi.insert((int)s.size());
            ans++;
        } else {
            mi.insert(s[0] - '0');
        }
    }
    return mi;
};

\end{lstlisting}

\subsection{random.cpp}
\begin{lstlisting}
#include <random>
#include <chrono>

mt19937_64 rng(chrono::steady_clock::now().time_since_epoch().count());
uniform_int_distribution<long long> dist(0, n);//valores min e max

int x = dist(rng);
\end{lstlisting}

\subsection{debug.cpp}
\begin{lstlisting}
//Debug

/*mais simples de todos, imprime desformatado na linha*/
void dbg_out() { cerr << endl; }
template<typename H, typename... T> 
void dbg_out(H h, T... t) { cerr << ' ' << h; dbg_out(t...); }

#define DEBUG

#if defined(DEBUG)
    #define winton (void)0
    #define debug(...) cerr << #__VA_ARGS__ << ':'; dbg_out(__VA_ARGS__);
#else
    #define winton ios_base::sync_with_stdio(false), cin.tie(NULL)
    #define debug(...) (void)0
#endif
//============================//


/*debug com cor (longo mas mt forte, imprime qualquer container tirando coisa com .pop() -> adicionar depois)*/
namespace dbg {
    const char* const RESET     = "\033[0m";
    const char* const BOLD_BLUE = "\033[1;34m";
    const char* const YELLOW    = "\033[33m";
    const char* const BOLD_WHITE= "\033[1;37m";

    template<typename T1, typename T2>
    ostream& operator<<(ostream& os, const pair<T1, T2>& p) { return os << '{' << p.first << ", " << p.second << '}'; }

    template<typename T_container, typename T = typename enable_if<!is_same_v<T_container, string> && !is_same_v<T_container, string_view>, typename T_container::value_type>::type>
    ostream& operator<<(ostream& os, const T_container& v) {
        os << '{';
        bool first = true;
        for (const T& x : v) { os << (first ? "" : ", ") << x, first = false; }
        return os << '}';
    }

    void debug_out(string_view) { cerr << endl; }
    template<typename H, typename... T>
    void debug_out(string_view s, H h, T... t) {
        auto cpos = s.find(',');
        cerr << YELLOW << s.substr(0, cpos) << RESET << " = ";
        cerr << BOLD_WHITE << h << RESET;
        if constexpr (sizeof...(t) > 0) {
            cerr << ", ";
            auto nx = s.find_first_not_of(" \t\n\r", cpos + 1);
            debug_out(s.substr(nx), t...);
        } else {
            cerr << endl;
        }
    }
} 
using namespace dbg;

#define DEBUG

#if defined(DEBUG)
    #define winton (void)0
    #define debug(...) cerr << BOLD_BLUE << "[" << __func__ << ":" << __LINE__ << "]" << RESET << " "; debug_out(#__VA_ARGS__, __VA_ARGS__)
#else
    #define winton ios_base::sync_with_stdio(false),cin.tie(NULL),cout.tie(NULL)
    #define debug(...) (void)0
#endif
//==========================================//


/*debug mais simples, so imprime vars mas imprime organizado na linha, menor e mais facil de codar*/
void debug_out(string s, int line) {cerr << endl; }
template<typename H, typename... T>
void debug_out(string s, int line, H h, T... t){
    do{
        cerr << line << ": " << s[0]; s = s.substr(1);
    }
    while (sz(s) and s[0] != ',');
    cerr << " = " << h;
    debug_out(s, line, t...);
}

#define DEBUG

#if defined(DEBUG)
    #define winton (void)0
    #define debug(...) debug_out(#__VA_ARGS__, __LINE__, __VA_ARGS__)
#else
    #define winton ios_base::sync_with_stdio(false), cin.tie(NULL)
    #define debug(...) (void)0
#endif
//==================================//


/*debug foda mas sem cor*/
namespace dbg {
    template<typename T1, typename T2>
    ostream& operator<<(ostream& os, const pair<T1, T2>& p) { return os << '{' << p.first << ", " << p.second << '}'; }

    template<typename T_container, typename T = typename enable_if<!is_same_v<T_container, string> && !is_same_v<T_container, string_view>, typename T_container::value_type>::type>
    ostream& operator<<(ostream& os, const T_container& v) {
        os << '{';
        bool first = true;
        for (const T& x : v) { os << (first ? "" : ", ") << x, first = false; }
        return os << '}';
    }

    void debug_out(string_view) { cerr << endl; }
    template<typename H, typename... T>
    void debug_out(string_view s, H h, T... t) {
        auto cpos = s.find(',');
        cerr << s.substr(0, cpos) << " = " << h;
        if constexpr (sizeof...(t) > 0) {
            cerr << ", ";
            auto nx = s.find_first_not_of(" \t\n\r", cpos + 1);
            debug_out(s.substr(nx), t...);
        } else {
            cerr << endl;
        }
    }
} 
using namespace dbg;

#define DEBUG

#if defined(DEBUG)
    #define winton (void)0
    #define debug(...) cerr << "[" << __func__ << ":" << __LINE__ << "]" << " "; debug_out(#__VA_ARGS__, __VA_ARGS__)
#else
    #define winton ios_base::sync_with_stdio(false),cin.tie(NULL),cout.tie(NULL)
    #define debug(...) (void)0
#endif
//====================================//


/*generic printer de todas as estruturas*/
// Generic printer for vector<T>
template <typename T>
ostream &operator<<(ostream &os, const vector<T> &v) {
    for (size_t i = 0; i < v.size(); ++i) {
        os << v[i];
        if (i + 1 < v.size()) os << ' ';
    }
    return os;
}

// Specialized printer for vector<vector<T>> (matrix)
template <typename T>
ostream &operator<<(ostream &os, const vector<vector<T>> &matrix) {
    for (size_t i = 0; i < matrix.size(); ++i) {
        os << matrix[i];
        if (i + 1 < matrix.size()) os << '\n';
    }
    return os;
}


//-------------PAIR-------------//
template<class A, class B> ostream& operator<<(ostream& os, const pair<A, B>& p) { 
    os << "(" << p.first << "," << p.second << ")"; 
    return os; 
}
//-------------VECTOR-------------//
template<typename T> ostream& operator<<(ostream& os, const vector<T>& vec) {
    os << "[ ";
    for(const auto& elem : vec) {
        os << elem << " ";
    }
    os << "]";
    return os;
}
//-------------SET-------------//
template<typename T> ostream& operator<<(ostream& os, const set<T>& s) {
    os << "{ ";
    for(const auto& elem : s) {
        os << elem << " ";
    }
    os << "}";
    return os;
}
//-------------MULTISET-------------//
template<typename T> ostream& operator<<(ostream& os, const multiset<T>& s) {
    os << "{ ";
    for(const auto& elem : s) {
        os << elem << " ";
    }
    os << "}";
    return os;
}
//-------------QUEUE-------------//
template<typename T> ostream& operator<<(ostream& os, queue<T> q) {
    // Print each element in the queue
    os << "{ ";
    while (!q.empty()) {
        os << q.front() << " ";
        q.pop();
    }
    os << "}";
    // Print a newline at the end
    return os;
}
//-------------DEQUE-------------//
template<typename T> ostream& operator<<(ostream& os, deque<T> q) {
    // Print each element in the queue
    os << "{ ";
    while (!q.empty()) {
        os << q.front() << " ";
        q.pop_front();
    }
    os << "}";
    // Print a newline at the end
    return os;
}
//-------------STACK-------------//
template<typename T> ostream& operator<<(ostream& os, stack<T> q) {
    // Print each element in the queue
    os << "{ ";
    while (!q.empty()) {
        os << q.top() << " ";
        q.pop();
    }
    os << "}";
    // Print a newline at the end
    return os;
}
//-------------PRIORITY_QUEUE-------------//
template<typename T> ostream& operator<<(ostream& os, priority_queue<T> q) {
    // Print each element in the queue
    os << "{ ";
    while (!q.empty()) {
        os << q.top() << " ";
        q.pop();
    }
    os << "}";
    // Print a newline at the end
    return os;
}
//-------------MAP-------------//
template<typename K, typename V> ostream& operator<<(ostream& os, const map<K, V>& m) {
    os << "{ ";
    for(const auto& pair : m) {
        os << pair.first << " : " << pair.second << ", ";
    }
    os << "}";
    return os;
}
 
template<typename T>
using min_pq = priority_queue<T, vector<T>, greater<T>>;
template<typename T> ostream& operator<<(ostream& os, min_pq<T> q) {
    // Print each element in the queue
    os << "{ ";
    while (!q.empty()) {
        os << q.top() << " ";
        q.pop();
    }
    os << "}";
    // Print a newline at the end
    return os;
}
\end{lstlisting}

\subsection{Template.cpp}
\begin{lstlisting}
//Template
#include <bits/stdc++.h>
using namespace std;

#define all(x) x.begin(), x.end()
#define rall(x) x.rbegin(), x.rend()
#define endl '\n'
#define int long long
#define ld long double

void dbg_out() { cerr << endl; }
template<typename H, typename... T> 
void dbg_out(H h, T... t) { cerr << ' ' << h; dbg_out(t...); }

#define DEBUG

#if defined(DEBUG)
    #define bg3 (void)0
    #define debug(...) cerr << #__VA_ARGS__ << ':'; dbg_out(__VA_ARGS__) ;
#else
    #define bg3 ios_base::sync_with_stdio(false), cin.tie(NULL)
    #define debug(...) (void)0
#endif

signed main(){
    bg3;
}
\end{lstlisting}

\end{document}
